% =============================================================================
%  Rapport technique — Plateforme CNC (Projets 28 & 15)
%  Compilation : latexmk -pdf rapport.tex   (ou pdflatex x2)
% =============================================================================
\documentclass[11pt,a4paper]{article}

\usepackage[T1]{fontenc}
\usepackage[utf8]{inputenc}
\usepackage{lmodern}
\usepackage[french]{babel}
\usepackage[a4paper,margin=2.2cm,headheight=14pt]{geometry}
\usepackage{microtype}
\usepackage{amsmath,amssymb}
\usepackage{graphicx}
\usepackage{float}
\usepackage[font=small,labelfont=bf]{caption}
\usepackage{booktabs,tabularx,longtable,array,multirow}
\usepackage{enumitem}
\usepackage[dvipsnames,table]{xcolor}
\usepackage{tikz}
\usetikzlibrary{arrows.meta,positioning,shapes.geometric,shapes.multipart,fit,backgrounds,calc,decorations.pathreplacing}
\usepackage[most]{tcolorbox}
\usepackage{listings}
\usepackage{fancyhdr}
\usepackage{titlesec}
\usepackage[hidelinks]{hyperref}

% ---------------------------------------------------------------- Couleurs CDC
\definecolor{navy}{HTML}{14365D}
\definecolor{royal}{HTML}{2B6CB0}
\definecolor{warm}{HTML}{C05621}
\definecolor{okgreen}{HTML}{2F855A}
\definecolor{todo}{HTML}{B7791F}
\definecolor{soft}{HTML}{EDF2F7}
\definecolor{codebg}{HTML}{F7FAFC}

% ---------------------------------------------------------------- Mise en page
\pagestyle{fancy}
\fancyhf{}
\fancyhead[L]{\small\color{navy}\textbf{Plateforme CNC} — Rapport technique}
\fancyhead[R]{\small\color{royal}Frontend \& Backend}
\fancyfoot[C]{\small\thepage}
\renewcommand{\headrulewidth}{0.4pt}

\titleformat{\section}{\Large\bfseries\color{navy}}{\thesection}{0.8em}{}
\titleformat{\subsection}{\large\bfseries\color{royal}}{\thesubsection}{0.6em}{}
\titleformat{\subsubsection}{\normalsize\bfseries\color{navy}}{\thesubsubsection}{0.5em}{}
\setlist{itemsep=2pt,topsep=4pt}
\setlength{\parindent}{0pt}
\setlength{\parskip}{5pt}
\sloppy
\setlength{\emergencystretch}{3em}

% ---------------------------------------------------------------- Commandes
\newcommand{\code}[1]{\texttt{\detokenize{#1}}}
\newcommand{\fait}{\textcolor{okgreen}{\textbf{Fait}}}
\newcommand{\afaire}{\textcolor{todo}{\textbf{À faire}}}
\newcommand{\partiel}{\textcolor{royal}{\textbf{Partiel}}}
\newcommand{\capture}[3][0.92]{%
  \begin{figure}[H]\centering
  \fcolorbox{black!15}{white}{\includegraphics[width=#1\linewidth]{images/#2}}
  \caption{#3}\end{figure}}

\newtcolorbox{note}[1][Remarque]{colback=soft,colframe=royal,fonttitle=\bfseries,title=#1,breakable,arc=2pt,boxrule=0.6pt}
\newtcolorbox{alerte}[1][Attention]{colback=orange!6,colframe=warm,fonttitle=\bfseries,title=#1,breakable,arc=2pt,boxrule=0.6pt}
\newtcolorbox{ok}[1][Résultat]{colback=green!4,colframe=okgreen,fonttitle=\bfseries,title=#1,breakable,arc=2pt,boxrule=0.6pt}

% ---------------------------------------------------------------- Code
\lstdefinelanguage{TypeScript}{
  keywords={const,let,var,function,return,if,else,for,of,in,new,throw,try,catch,async,await,
    import,export,from,type,interface,class,extends,readonly,true,false,null,undefined,typeof,as,default},
  sensitive=true, comment=[l]{//}, morecomment=[s]{/*}{*/},
  morestring=[b]', morestring=[b]", morestring=[b]`}
\lstdefinelanguage{Prisma}{
  keywords={model,enum,generator,datasource,String,Int,Float,Boolean,DateTime,Json},
  sensitive=true, comment=[l]{//}, morestring=[b]"}
\lstset{
  basicstyle=\ttfamily\footnotesize, keywordstyle=\color{royal}\bfseries,
  commentstyle=\color{black!50}\itshape, stringstyle=\color{warm},
  backgroundcolor=\color{codebg}, frame=single, rulecolor=\color{black!15},
  breaklines=true, columns=fullflexible, keepspaces=true, showstringspaces=false,
  numbers=left, numberstyle=\tiny\color{black!40}, xleftmargin=1.5em, framexleftmargin=1.5em,
  literate={é}{{\'e}}1 {è}{{\`e}}1 {à}{{\`a}}1 {ê}{{\^e}}1 {ç}{{\c c}}1 {ô}{{\^o}}1
           {É}{{\'E}}1 {î}{{\^i}}1 {ù}{{\`u}}1 {→}{{$\rightarrow$}}1
           {û}{{\^u}}1 {â}{{\^a}}1 {ï}{{\"i}}1 {ë}{{\"e}}1 {ü}{{\"u}}1 {À}{{\`A}}1
           {È}{{\`E}}1 {Ê}{{\^E}}1 {Ç}{{\c C}}1 {œ}{{\oe}}1 {«}{{\guillemotleft}}1 {»}{{\guillemotright}}1
}

% ---------------------------------------------------------------- Styles TikZ
\tikzset{
  box/.style={draw=navy,rounded corners=3pt,fill=white,align=center,font=\small,inner sep=5pt,line width=0.6pt},
  svc/.style={box,fill=royal!8,minimum width=3.2cm,minimum height=1.1cm},
  done/.style={box,fill=green!10,draw=okgreen},
  skel/.style={box,fill=orange!8,draw=warm,dashed},
  db/.style={cylinder,shape border rotate=90,aspect=0.18,draw=navy,fill=navy!8,align=center,font=\small,minimum height=1.3cm,minimum width=2.2cm},
  arr/.style={-{Stealth[length=2.2mm]},line width=0.6pt,color=navy},
  darr/.style={arr,dashed},
  lbl/.style={font=\scriptsize,fill=white,inner sep=1pt},
}

\hypersetup{pdftitle={Plateforme CNC - Rapport technique},pdfauthor={A. Lemrani, M. Rouijel}}

% =============================================================================
\begin{document}

\begin{titlepage}
  \centering
  {\color{navy}\rule{\linewidth}{1.2pt}}\\[1.2em]
  {\large\textsc{Royaume du Maroc — CPGE \& Cycle Ingénieur}}\\[2.2em]
  {\Huge\bfseries\color{navy} Rapport technique du projet}\\[0.6em]
  {\LARGE\color{royal} Plateforme de préparation au CNC}\\[0.4em]
  {\large Projet 28 \& extension Projet 15 (QCM adaptatif)}\\[1.2em]
  {\color{navy}\rule{\linewidth}{1.2pt}}\\[2.5em]

  \begin{tikzpicture}[node distance=6mm]
    \node[box,fill=royal!10,minimum width=4.2cm,minimum height=1.2cm] (f) {\textbf{frontend/}\\ Next.js 16 · React 19};
    \node[box,fill=green!10,draw=okgreen,minimum width=4.2cm,minimum height=1.2cm,right=2.8cm of f] (b) {\textbf{backend/}\\ Node.js · 7 processus};
    \draw[arr,{Stealth[length=2.2mm]}-{Stealth[length=2.2mm]}] (f) -- node[above=2pt,font=\scriptsize,text=navy]{HTTP / JSON} (b);
  \end{tikzpicture}\\[3em]

  \begin{minipage}{0.8\linewidth}\centering\large
    Ce document explique \textbf{ce qui a été fait}, \textbf{ce qu'il reste à faire} et \textbf{la logique du code} du frontend et du backend, avec des schémas et des captures d'écran de l'application.
  \end{minipage}\\[\fill]

  \begin{tabular}{ll}
    \textbf{Réalisé par :} & Achraf Lemrani, Mohammed Rouijel\\
    \textbf{Encadré par :} & M. Hosni\\
    \textbf{Référence :} & Cahier des charges (\code{docs/cdc.pdf})\\
    \textbf{Date :} & \today\\
  \end{tabular}
\end{titlepage}

\tableofcontents
\newpage

\section{Introduction}

\subsection{Rappel du besoin (cahier des charges)}

Le cahier des charges (CDC) décrit une plateforme web de préparation au \textbf{Concours National Commun} (CNC) pour les étudiants de CPGE des filières \textbf{MP}, \textbf{PSI} et \textbf{TSI}, sur quatre matières : Mathématiques, Physique, Sciences de l'Ingénieur et Informatique. Les modules fonctionnels sont :

\begin{enumerate}[label=\textbf{M\arabic*}]
  \item Authentification, profil étudiant et navigation libre Sup / Spé (JWT + bcrypt) ;
  \item Cours par leçon et exercices en double mode (PDF classique et mode interactif type Bibm@th) ;
  \item Problèmes, annales CNC 2010–2026, simulateur chronométré et « Mes Problèmes Rédigés » ;
  \item Bibliothèque d'ouvrages de référence ;
  \item Dashboard personnel : Streak, heatmap, statistiques par matière ;
  \item Contact admin et signalement d'errata (tickets) ;
  \item Fonctionnalités optionnelles (fiches express, TIPE, simulateur de rang, oraux, carnet d'erreurs).
\end{enumerate}

L'extension \textbf{Projet 15} ajoute un QCM intelligent : indice de maîtrise $M(u,c)$, difficulté adaptative et répétition espacée.

Sur le plan technique, le CDC (section 3.1) impose une \textbf{architecture microservices en Node.js (TypeScript)} : une API Gateway et six microservices, PostgreSQL 16 avec Prisma (un schéma logique par service), Redis et un stockage objet S3/MinIO.

\subsection{Objectif de ce document}

Ce rapport sert à relire et comprendre le projet. Il répond à trois questions :

\begin{description}[leftmargin=1.5em,style=nextline]
  \item[Ce qui a été fait] l'organisation du dépôt en deux applications séparées, le frontend (prototype complet) et le backend (architecture complète, Sprint 1 implémenté et testé).
  \item[Ce qu'il reste à faire] les sprints 2 à 5 du backend, le branchement du frontend sur l'API et les écarts au CDC.
  \item[La logique du code] le cheminement d'une requête, le rôle de chaque fichier et les algorithmes (JWT, rotation des refresh tokens, Streak, QCM adaptatif).
\end{description}

\subsection{Synthèse de l'état d'avancement}

\begin{table}[H]\centering\small
\begin{tabularx}{\linewidth}{>{\bfseries}l X c}
\toprule
Élément & Contenu & État\\
\midrule
Organisation du dépôt & \code{frontend/}, \code{backend/}, \code{docs/}, \code{archive/} séparés & \fait\\
Frontend (prototype) & 7 routes, 42 composants, environ 37\,700 lignes, données de démonstration & \fait\\
API Gateway & Routage, JWT, rate limiting, CORS, santé agrégée & \fait\\
\code{auth-user-service} & Inscription, connexion, refresh avec rotation, déconnexion, profil, liste admin & \fait\\
Package \code{@cnc/common} & Configuration, erreurs, contexte d'identité, application type & \fait\\
5 autres microservices & Schéma Prisma complet, routes déclarées (réponse 501) & \partiel\\
Base de données & PostgreSQL 16 : 6 schémas, 23 tables créées ; migrations versionnées à mettre en place (section~\ref{sec:bdd}) & \partiel\\
Moteur de Streak, moteur QCM & Formules du CDC codées et testées (logique pure) & \fait\\
Logique métier des sprints 2 à 5 & Catalogue, examens, analytique, contact, QCM & \afaire\\
Branchement frontend $\leftrightarrow$ API & Remplacer les données en dur par des appels HTTP & \afaire\\
\bottomrule
\end{tabularx}
\caption{État d'avancement global}
\end{table}

\section{Organisation du dépôt}

\subsection{Avant et après la réorganisation}

Avant la réorganisation, le frontend Next.js occupait la racine du dossier \code{projet/}. Il y côtoyait un dossier \code{backend/} incomplet : son \code{package.json} appelait \code{gateway.ts} et \code{auth.ts}, deux fichiers qui n'existaient pas. À la même racine se trouvaient aussi le CDC, des captures, des archives zip, un dossier \code{tmp/} de 580\,Mo et une jonction Windows (\code{test_junc}) qui pointait vers \code{app/}.

\begin{figure}[H]\centering
\begin{tikzpicture}[
  folder/.style={box,minimum width=2.6cm,font=\footnotesize\ttfamily,inner sep=3pt},
  every node/.append style={font=\footnotesize}]
  % AVANT
  \node[font=\bfseries\color{warm}] at (0,0.9) {AVANT};
  \node[folder,fill=red!5,draw=warm,minimum width=4.2cm,align=left,anchor=north] (av) at (0,0.5) {%
    projet/\\
    \ app/ components/ lib/ public/\\
    \ scripts/ node\_modules/ .next/\\
    \ package.json next.config.mjs ...\\
    \ backend/ \textcolor{warm}{(incomplet)}\\
    \ cdc.pdf audit-ui-cdc.md\\
    \ tmp/ output/ *.zip *.png\\
    \ test\_junc \textcolor{warm}{(jonction)}};
  \draw[arr,line width=1.2pt] (2.4,-1.3) -- (4.2,-1.3);
  % APRES
  \node[font=\bfseries\color{okgreen}] at (8.2,0.9) {APRÈS};
  \node[folder,fill=royal!8] (fe) at (6.3,0) {frontend/};
  \node[folder,fill=green!10,draw=okgreen] (be) at (10.1,0) {backend/};
  \node[folder,fill=soft] (dc) at (6.3,-1.6) {docs/};
  \node[folder,fill=black!5,draw=black!40] (ar) at (10.1,-1.6) {archive/};
  \node[font=\scriptsize,align=center,below=1mm of fe] {Next.js autonome};
  \node[font=\scriptsize,align=center,below=1mm of be] {microservices Node.js};
  \node[font=\scriptsize,align=center,below=1mm of dc] {CDC, audit, rapport};
  \node[font=\scriptsize,align=center,below=1mm of ar] {maquettes, tmp, zips};
  \node[folder,fill=white] (rd) at (8.2,-3.1) {README.md};
\end{tikzpicture}
\caption{Réorganisation du dépôt en deux applications indépendantes}
\end{figure}

\subsection{Arborescence finale}

\begin{lstlisting}[numbers=none,language={},caption={Arborescence de premier niveau}]
projet/
+-- README.md                 Explique l'organisation et les commandes
+-- frontend/                 Application Next.js (package.json et node_modules propres)
|   +-- app/                  Routes App Router : /, /login, /signup, /admin, /pochettes, /parametres
|   +-- components/           42 composants React (dont document-studio/ et ui/)
|   +-- lib/                  auth.ts (session), bases d'annales, manifestes JSON
|   +-- public/               579 PDF (annales, problèmes), pdf.js, logos
|   +-- scripts/              téléchargement d'annales, génération de PDF, logos
+-- backend/                  Monorepo npm workspaces (package.json racine)
|   +-- gateway/              API Gateway (port 4000)
|   +-- packages/common/      code partagé par tous les services
|   +-- services/             6 microservices (ports 4001 à 4006)
|   +-- compose.yaml          PostgreSQL 16, Redis 7, MinIO
|   +-- .env.example          variables d'environnement documentées
+-- docs/                     cdc.pdf, audit-ui-cdc.md, test-accounts.md, rapport/
+-- archive/                  anciens fichiers, conservés mais non utilisés
\end{lstlisting}

\begin{note}[Ce qui garantit l'indépendance des deux applications]
\begin{itemize}
  \item Chaque application a ses propres \code{package.json}, \code{node_modules} et \code{tsconfig}. On peut travailler sur l'une sans installer l'autre.
  \item Elles communiquent \emph{uniquement} en HTTP/JSON, à travers la gateway (\code{http://localhost:4000/api/...}).
  \item Après le déplacement, le build de production du frontend (\code{next build}) réussit, et les scripts du frontend se repèrent par rapport à leur propre dossier (\code{import.meta.url}).
\end{itemize}
\end{note}

\begin{alerte}[Points d'attention sur la machine de développement]
\begin{itemize}
  \item Le port \textbf{5433} était déjà pris par un autre conteneur (\code{launchops-db}). PostgreSQL du projet utilise donc le port \textbf{5434}.
  \item Au démarrage, Docker Desktop relance un conteneur \code{open-webui} sur le port \textbf{3000}. Si ce conteneur tourne, \code{next dev} se rabattra sur le port 3001.
  \item Le dossier \code{archive/tmp/github-publish} contient un ancien clone git. Il a été déplacé tel quel.
\end{itemize}
\end{alerte}

\section{Architecture globale}

\subsection{Vue d'ensemble}

L'architecture reprend la section 3.1 du CDC. Le navigateur ne connaît qu'une seule adresse, celle de l'\textbf{API Gateway}. C'est la gateway qui vérifie l'identité de l'utilisateur, puis transmet la requête au microservice concerné.

\begin{figure}[H]\centering
\begin{tikzpicture}[node distance=5mm]
  % Client
  \node[box,fill=royal!12,minimum width=12.6cm,minimum height=1cm] (fe) {\textbf{Frontend} — Next.js 16 / React 19 / Tailwind / KaTeX / pdf.js \hfill \texttt{:3000}};
  % Gateway
  \node[box,fill=navy!10,minimum width=12.6cm,minimum height=1.25cm,below=9mm of fe] (gw)
    {\textbf{API Gateway} \texttt{:4000} — CORS · rate limiting · vérification JWT · routage \texttt{/api/*}};
  \draw[arr] (fe) -- node[lbl,right]{HTTPS / JSON + \texttt{Authorization: Bearer}} (gw);

  % Services
  \begin{scope}[node distance=2mm]
  \node[done,minimum width=1.95cm,minimum height=1.5cm,below=11mm of gw.south west,anchor=north west,font=\scriptsize] (s1) {\textbf{auth-user}\\ \texttt{:4001}\\ Sprint 1};
  \node[skel,minimum width=1.95cm,minimum height=1.5cm,right=of s1,font=\scriptsize] (s2) {\textbf{catalog-}\\\textbf{course} \texttt{:4002}\\ Sprint 2};
  \node[skel,minimum width=1.95cm,minimum height=1.5cm,right=of s2,font=\scriptsize] (s3) {\textbf{exam-}\\\textbf{submission}\\ \texttt{:4003} S3};
  \node[skel,minimum width=1.95cm,minimum height=1.5cm,right=of s3,font=\scriptsize] (s4) {\textbf{analytics-}\\\textbf{streak}\\ \texttt{:4004} S4};
  \node[skel,minimum width=1.95cm,minimum height=1.5cm,right=of s4,font=\scriptsize] (s5) {\textbf{admin-}\\\textbf{contact}\\ \texttt{:4005} S4};
  \node[skel,minimum width=1.95cm,minimum height=1.5cm,right=of s5,font=\scriptsize] (s6) {\textbf{adaptive-}\\\textbf{qcm}\\ \texttt{:4006} S5};
  \end{scope}
  \foreach \s/\p in {s1/auth,s2/catalog,s3/exams,s4/analytics,s5/contact,s6/qcm}{
    \draw[arr] (gw.south -| \s.north) -- node[lbl,font=\tiny]{\texttt{/\p}} (\s.north);
  }

  % Common
  \node[box,fill=soft,minimum width=12.6cm,below=4mm of s1.south west,anchor=north west,font=\scriptsize] (cm)
    {\texttt{packages/common} : configuration, erreurs, \texttt{requireGateway}, \texttt{requireUser}, \texttt{createServiceApp}};

  % Data
  \node[db,below=8mm of cm.south,xshift=-3.6cm,minimum width=5.2cm] (pg) {\textbf{PostgreSQL 16} \texttt{:5434}\\[-1pt]\scriptsize schémas : auth · catalog · exam\\[-2pt]\scriptsize analytics · contact · qcm};
  \node[db,right=8mm of pg,fill=red!6] (rd) {\textbf{Redis}\\\scriptsize cache, événements};
  \node[db,right=8mm of rd,fill=orange!8] (mn) {\textbf{MinIO / S3}\\\scriptsize PDF, copies};
  \draw[arr] (cm.south -| pg.north) -- node[lbl,font=\scriptsize]{Prisma} (pg.north);
  \draw[darr] (cm.south -| rd.north) -- (rd.north);
  \draw[darr] (cm.south -| mn.north) -- (mn.north);

  % Legend
  \node[done,font=\scriptsize,below=6mm of pg.south west,anchor=north west] (l1) {implémenté};
  \node[skel,font=\scriptsize,right=3mm of l1] (l2) {squelette (schéma + routes 501)};
  \draw[darr] ($(l2.east)+(4mm,0)$) -- ++(9mm,0) node[right,font=\scriptsize,text=black]{prévu (Redis, MinIO pas encore utilisés)};
\end{tikzpicture}
\caption{Architecture microservices : en vert ce qui est implémenté, en orange les squelettes}
\label{fig:archi}
\end{figure}

\subsection{Principes de conception}

\begin{description}[leftmargin=1.5em,style=nextline]
  \item[Un seul point d'entrée] Le frontend appelle uniquement \code{/api/...} sur la gateway. Les services ne sont jamais appelés directement par le navigateur.
  \item[JWT vérifié une seule fois] La gateway valide le jeton, puis transmet l'identité au service dans les en-têtes \code{x-user-id} et \code{x-user-role}. Les services n'ont pas besoin du secret JWT ; seul \code{auth-user-service} le connaît, car c'est lui qui signe les jetons.
  \item[Confiance gateway $\rightarrow$ service] Chaque requête transmise porte un secret partagé, \code{x-internal-secret}. Un service refuse (403) toute requête qui ne le présente pas : un client qui l'appellerait directement ne peut donc pas usurper une identité.
  \item[Une base par service] Il n'y a qu'un seul serveur PostgreSQL, mais chaque service a son propre schéma et son propre client Prisma (généré dans \code{src/generated/prisma}). Aucun service ne lit les tables d'un autre. Les références entre services, comme \code{userId}, sont de simples UUID sans clé étrangère.
  \item[Code commun factorisé] Le paquet \code{@cnc/common} évite la duplication entre services : chargement du \code{.env}, format des erreurs, contrôle d'accès et construction de l'application Express.
\end{description}

\subsection{Cycle de vie d'une requête authentifiée}

\begin{figure}[H]\centering
\begin{tikzpicture}[x=1cm,y=-0.62cm]
  \foreach \x/\n/\c in {0/Navigateur/royal,4.2/Gateway/navy,8.4/auth-user-service/okgreen,12.4/PostgreSQL/navy}{
    \node[box,draw=\c,fill=\c!10,minimum width=2.9cm] at (\x,0) {\textbf{\n}};
    \draw[black!30,dashed] (\x,0.7) -- (\x,13.3);
  }
  \draw[arr] (0,1.5) -- node[lbl,above]{\texttt{GET /api/users/me} + Bearer} (4.2,1.5);
  \node[box,fill=soft,font=\scriptsize,align=left,text width=3.3cm] at (4.2,3.7) {1. supprime \texttt{x-user-*} reçus\\2. \texttt{jwt.verify(token)}\\3. ajoute \texttt{x-user-id}, \texttt{x-user-role}, \texttt{x-internal-secret}};
  \draw[arr] (4.2,6) -- node[lbl,above]{\texttt{GET /users/me} (préfixe \texttt{/api} retiré)} (8.4,6);
  \node[box,fill=soft,font=\scriptsize,align=left,text width=3.4cm] at (8.4,8.1) {4. \texttt{requireGateway}\\5. \texttt{requireUser()}\\6. \texttt{users.getUser(id)}};
  \draw[arr] (8.4,10.2) -- node[lbl,above]{\texttt{findUnique}} (12.4,10.2);
  \draw[darr] (12.4,11) -- (8.4,11);
  \draw[darr] (8.4,12) -- node[lbl,above]{200 + JSON (sans \texttt{passwordHash})} (0,12);
\end{tikzpicture}
\caption{Diagramme de séquence : lecture du profil de l'utilisateur connecté}
\end{figure}

\section{Frontend : ce qui a été fait}

\subsection{Pile technique et volume}

\begin{table}[H]\centering\small
\begin{tabularx}{\linewidth}{l X}
\toprule
\textbf{Bibliothèque} & \textbf{Rôle dans le projet}\\
\midrule
Next.js 16 (App Router), React 19, TypeScript & Structure des pages ; tous les écrans sont des composants client (\code{'use client'})\\
Tailwind CSS v4, \code{tw-animate-css} & Styles ; mode sombre piloté par la classe \code{.dark} sur \code{<html>}\\
KaTeX 0.19 & Rendu des formules (\code{MathFormula.tsx} appelle \code{katex.renderToString})\\
\code{pdfjs-dist} 6.5 & Liseuse PDF intégrée (\code{document-studio/}) : pages, recherche, signets, mode examen\\
\code{framer-motion} & Animations (cartes de dossiers, lecteur d'épreuve)\\
\code{lucide-react}, \code{@phosphor-icons/react} & Icônes\\
\bottomrule
\end{tabularx}
\caption{Bibliothèques principales du frontend}
\end{table}

\begin{table}[H]\centering\small
\begin{tabular}{lrr}
\toprule
\textbf{Dossier} & \textbf{Fichiers} & \textbf{Lignes}\\
\midrule
\code{app/} & 9 & 2\,313\\
\code{components/} & 42 & 25\,471\\
\code{lib/} (dont 2 manifestes JSON) & 8 & 9\,326\\
\code{scripts/} & 6 & 663\\
\midrule
\textbf{Total} & 65 & \textbf{37\,773}\\
\bottomrule
\end{tabular}
\quad
\begin{tabular}{lr}
\toprule
\textbf{Ressources} & \textbf{Volume}\\
\midrule
PDF d'annales & 559 (568\,Mo)\\
PDF de problèmes & 20 (1,9\,Mo)\\
Épreuves référencées & 437\\
Chapitres de cours & 26\\
\bottomrule
\end{tabular}
\caption{Volume de code et de contenu du frontend}
\end{table}

\subsection{Routes et navigation}

\begin{table}[H]\centering\small
\begin{tabularx}{\linewidth}{l l X}
\toprule
\textbf{Route} & \textbf{Fichier} & \textbf{Contenu}\\
\midrule
\code{/} & \code{app/page.tsx} & Application étudiant : 6 onglets (Dashboard, Cours, Annales, Problèmes, QCM, Contact)\\
\code{/login} & \code{app/login/page.tsx} & \code{LoginPage}, puis redirection vers \code{/admin} ou \code{/} selon le rôle\\
\code{/signup} & \code{app/signup/page.tsx} & \code{SignUpPage} (nom, e-mail, centre, année, filière)\\
\code{/admin} & \code{app/admin/page.tsx} & \code{AdminBackOffice} : messages, catalogue, annales, QCM\\
\code{/pochettes} & \code{app/pochettes/page.tsx} & Espace personnel en dossiers, éditeur de pages (LaTeX, PDF)\\
\code{/parametres} & \code{app/parametres/page.tsx} & Profil, cursus, sécurité, notifications\\
\bottomrule
\end{tabularx}
\caption{Routes de l'application}
\end{table}

\begin{figure}[H]\centering
\resizebox{\linewidth}{!}{\begin{tikzpicture}[
  pg/.style={box,fill=royal!8,font=\scriptsize,minimum height=0.8cm,text width=2.2cm},
  sub/.style={box,fill=white,font=\scriptsize,text width=2.2cm},
  node distance=3mm and 4mm]
  \node[box,fill=navy!12,font=\small\bfseries] (root) {\texttt{app/page.tsx} — \texttt{activeTab} + \texttt{?tab=}};
  \node[pg,below=9mm of root,xshift=-6.6cm] (d) {Dashboard\\\texttt{dashboard}};
  \node[pg,right=of d] (c) {Cours \& TD\\\texttt{cours}};
  \node[pg,right=of c] (a) {Annales\\\texttt{annales}};
  \node[pg,right=of a] (p) {Problèmes\\\texttt{problemes}};
  \node[pg,right=of p] (q) {QCM\\\texttt{qcm}};
  \node[pg,right=of q] (ct) {Contact\\\texttt{contact}};
  \foreach \n in {d,c,a,p,q,ct}{\draw[arr] (root.south) -- (\n.north);}
  \node[sub,below=of d] {StatCard, RevisionChart, ActivityBar, Achievement\-BadgesCard};
  \node[sub,below=of c] (c1) {Niveaux (Sup/Spé)};
  \node[sub,below=of c1] (c2) {Matières (dossiers)};
  \node[sub,below=of c2] (c3) {Chapitres};
  \node[sub,below=of c3] (c4) {Cours : KaTeX $\leftrightarrow$ planche PDF};
  \node[sub,below=of c4] (c5) {TD : indices, corrigé, auto-évaluation, QCM flash};
  \foreach \i/\j in {c1/c2,c2/c3,c3/c4,c4/c5}{\draw[arr] (\i) -- (\j);}
  \node[sub,below=of a] (a1) {5 concours};
  \node[sub,below=of a1] (a2) {Sessions (années)};
  \node[sub,below=of a2] (a3) {Épreuves};
  \node[sub,below=of a3] (a4) {Lecteur : normal ou concours chronométré};
  \foreach \i/\j in {a1/a2,a2/a3,a3/a4}{\draw[arr] (\i) -- (\j);}
  \node[sub,below=of p] (p1) {10 problèmes classiques};
  \node[sub,below=of p1] {Mode focus + lecteur PDF};
  \draw[arr] (p1) -- ++(0,-0.75);
  \node[sub,below=of q] (q1) {Choix matière};
  \node[sub,below=of q1] {Questions, chrono 90\,s, score};
  \draw[arr] (q1) -- ++(0,-0.75);
  \node[sub,below=of ct] {Tickets, messages, pièces jointes};
\end{tikzpicture}}
\caption{Arbre de navigation de l'espace étudiant}
\end{figure}

Les onglets Cours, Annales et QCM sont tous les trois rendus par le même composant, \code{CourseExerciseQcmDoubleMode.tsx} (5\,276 lignes), avec une valeur différente de la prop \code{initialSection}.

\paragraph{Logique de navigation.} \code{app/page.tsx} garde l'onglet actif dans \code{activeTab}. \code{navigateToTab()} met à jour cet état et écrit \code{?tab=<id>} dans l'URL avec \code{history.pushState}. Au chargement et sur \code{popstate}, \code{syncTabFromUrl()} relit le paramètre, en n'acceptant que les onglets d'une liste blanche. Les boutons Précédent et Suivant du navigateur fonctionnent donc sans rechargement.

\paragraph{Session (simulée).} \code{lib/auth.ts} enregistre un objet \code{UserProfile} dans \code{localStorage} (clé \code{learning_cnc_user_session}). Les pages le lisent avec \code{getStoredUser() || DEMO_STUDENT}. Elles se synchronisent grâce à l'événement \code{nexora-user-updated}, émis par la page des paramètres.

\paragraph{Thème.} \code{ThemeToggle} lit \code{localStorage['nexora-theme']}, à défaut \code{prefers-color-scheme}, puis bascule la classe \code{dark} sur \code{<html>}. Le fond animé WebGL (\code{SilkBackground}) observe cette classe avec un \code{MutationObserver}.

\subsection{Écrans réalisés}

\subsubsection{Authentification (M1)}
\capture[0.75]{01-login.png}{Connexion : comptes de démonstration étudiant et administrateur}
\capture[0.75]{02-signup.png}{Inscription : prénom, nom, e-mail, centre CPGE, année et filière (UC01)}

\textbf{Logique :} \code{LoginPage} renvoie \code{DEMO_ADMIN} si l'e-mail est celui de l'administrateur, sinon \code{DEMO_STUDENT}, puis appelle \code{saveUserSession()}. \textbf{Le mot de passe n'est pas vérifié} : c'est précisément le rôle du backend (section~\ref{sec:auth}).

\subsubsection{Dashboard (M5)}
\capture{03-dashboard.png}{Dashboard étudiant : progression, temps de révision par matière, objectifs, activité}
\capture{04-dashboard-dark.png}{Le même écran en mode sombre (exigence non fonctionnelle 7.1)}

Les valeurs affichées (72\,\%, 12 jours, 15 annales, 254\,h…) sont \textbf{écrites en dur}. La heatmap (\code{ActivityBar}) est générée par une formule pseudo-aléatoire. Ces données viendront de \code{analytics-streak-service}.

\subsubsection{Cours et exercices en double mode (M2)}
\capture{05-cours-niveaux.png}{Choix du niveau : accès libre Sup et Spé}
\capture{13-matieres.png}{Choix de la matière (cartes « dossiers »)}
\capture{14-chapitres.png}{Liste des chapitres d'une matière, avec accès à la partie Cours ou TD}
\capture{15-cours-interactif.png}{Partie Cours en mode « Fiche interactive KaTeX » : théorie, théorèmes, formules}
\capture{16-cours-pdf.png}{Partie Cours en mode « Planche PDF » (rendu HTML façon document officiel)}
\capture{18-td-indice-solution.png}{Partie TD : exercice référencé (\texttt{EXO-MP-MATH-204}), indications progressives, corrigé pas à pas et QCM flash}

\textbf{Logique :} toutes les données (26 chapitres, 28 exercices, 26 questions) sont dans le tableau \code{COURSES_DATABASE} de \code{CourseExerciseQcmDoubleMode.tsx}. Chaque exercice a le type suivant :
\begin{lstlisting}[language=TypeScript,numbers=none]
interface CourseExercise {
  id: string; refCode: string; source: string; difficultyStars: number
  title: string; statement: string; statementMath?: string[]
  hints: string[]                                  // indices progressifs
  solutionSteps: { stepTitle: string; content: string; math?: string }[]
}
\end{lstlisting}
Le composant mémorise \code{unlockedHints[exoId]}, le nombre d'indices déjà dévoilés, et \code{revealedSolutions[exoId]}. L'auto-évaluation (Maîtrisé, Avec aide, À revoir) est stockée dans \code{exerciseEvaluations}, \textbf{uniquement en mémoire}. Ce type correspond presque exactement au modèle \code{Exercise} du backend (\code{statementLatex}, \code{hintsLatex[]}, \code{solutionLatex}).

\subsubsection{Annales CNC et simulateur d'épreuve (M3)}
\capture{06-annales.png}{Portail des concours : 437 épreuves référencées, dont 156 CNC avec de vrais PDF}
\capture{19-annales-cnc.png}{Sessions d'annales CNC (2000 à 2025)}
\capture{19b-epreuves-2024.png}{Épreuves d'une session, filtres par matière, visionnage normal ou mode concours}
\capture{21-mode-concours.png}{Mode concours : phase de préparation, chronomètre, brouillon, remise de copie}

\textbf{Logique du chronomètre} (\code{ConcoursPdfExamViewer.tsx}) : l'épreuve passe par trois phases.
\begin{center}
\begin{tikzpicture}[node distance=26mm]
  \node[box,fill=yellow!15] (p) {\texttt{prep}\\\scriptsize 300\,s, énoncé masqué};
  \node[box,fill=royal!10,right=of p] (a) {\texttt{active}\\\scriptsize durée $\times$ 3600\,s};
  \node[box,fill=green!10,right=of a] (s) {\texttt{submitted}\\\scriptsize corrigé déverrouillé};
  \draw[arr] (p) -- node[above=2pt,font=\scriptsize,align=center]{chrono à 0\\ou « Débuter »} (a);
  \draw[arr] (a) -- node[above=2pt,font=\scriptsize,align=center]{chrono à 0\\ou « Rendre »} (s);
  \node[box,fill=red!6,draw=warm,font=\scriptsize,below=6mm of a] (i) {sortie de fenêtre (\texttt{blur}, \texttt{visibilitychange}) $\rightarrow$ infraction enregistrée};
  \draw[darr] (a) -- (i);
\end{tikzpicture}
\end{center}
Un \code{setInterval} d'une seconde décrémente le temps restant. Les annales sont décrites par deux manifestes JSON, générés par \code{scripts/download-*-annales.mjs}, qui pointent vers \code{public/annales/...}.

\subsubsection{Problèmes, QCM et contact}
\capture{07-problemes.png}{Banque de problèmes classiques (énoncé et corrigé PDF générés par \texttt{build-problem-pdfs.cjs})}
\capture{22-qcm-question.png}{QCM par matière : question, choix rendus en KaTeX, chronomètre de 90\,s}
\capture{09-contact.png}{Contact admin : tickets, messagerie, pièces jointes (M6)}

\subsubsection{Back-office et espace personnel}
\capture{10-admin.png}{Back-office administrateur : demandes à traiter, annales, statistiques}
\capture{11-pochettes.png}{Pochettes : espace personnel organisé en dossiers}
\capture{12-parametres.png}{Paramètres : profil, cursus, sécurité, notifications}

\subsection{Bilan du frontend}

\begin{ok}[Points forts]
L'interface est complète et cohérente avec les maquettes du CDC : double mode d'exercices, rendu KaTeX, liseuse PDF avec mode examen, mode sombre, mise en page responsive, 559 PDF d'annales réellement intégrés, et une navigation qui conserve l'onglet dans l'URL.
\end{ok}

\begin{alerte}[Limites : c'est un prototype]
\begin{itemize}
  \item \textbf{Aucun appel réseau} (aucun \code{fetch}) : toutes les données sont en dur ou dans \code{localStorage}, et rien n'est conservé côté serveur.
  \item \textbf{Pas de vraie sécurité} : le mot de passe n'est pas vérifié, la route \code{/admin} est accessible sans contrôle de rôle, et \code{clearUserSession()} n'est jamais appelée à la déconnexion.
  \item \textbf{Composants très volumineux} : \code{CourseExerciseQcmDoubleMode.tsx} fait 5\,276 lignes et \code{AdminBackOffice.tsx} 3\,726 lignes. Ils mélangent données et affichage.
  \item \textbf{Code orphelin} : \code{StudentDrawer}, \code{PersonalWorkspaceModal} et \code{glass-folder-cards} ne sont rendus nulle part ; l'onglet \code{fiches} est inaccessible.
  \item Le chronomètre du QCM flash est figé sur « 01:45 » ; la répétition espacée n'est pas reliée à l'interface.
  \item \code{next.config.mjs} contient \code{ignoreBuildErrors: true} : les erreurs TypeScript sont masquées au build.
\end{itemize}
\end{alerte}

\section{Backend : ce qui a été fait}

\subsection{Organisation en monorepo}

Le backend est un \textbf{monorepo npm workspaces}. Un seul \code{npm install} à la racine de \code{backend/} installe les dépendances de la gateway, du paquet commun et des six services. Chaque service reste pourtant un paquet indépendant, avec son \code{package.json}, son schéma Prisma et ses scripts.

\begin{lstlisting}[numbers=none,language={},caption={Arborescence du backend (1\,336 lignes de TypeScript et Prisma écrites à la main)}]
backend/
+-- package.json            workspaces + scripts globaux (dev, typecheck, test, db:*)
+-- tsconfig.base.json      options TypeScript communes (strict, ES2023, Bundler)
+-- compose.yaml            postgres:16 (5434), redis:7 (6379), minio (9000/9001)
+-- .env.example            toutes les variables, une URL de base par service
+-- packages/common/src/
|   +-- env.ts              charge backend/.env, requireEnv(), portEnv()
|   +-- errors.ts           HttpError, notFoundHandler, errorHandler (Zod -> 400)
|   +-- auth-context.ts     en-têtes x-user-*, requireGateway(), requireUser(), currentUser()
|   +-- service-app.ts      createServiceApp() : /health + sécurité + JSON + erreurs
|   +-- not-implemented.ts  notImplemented('Sprint N') -> 501
+-- gateway/src/
|   +-- index.ts            CORS, rate limit, /api/health, proxys
|   +-- authenticate.ts     vérification JWT + injection d'identité
|   +-- routes.ts           table préfixe -> service, routes publiques
+-- services/
    +-- auth-user-service/      IMPLÉMENTÉ (Sprint 1)
    |   +-- prisma/schema.prisma, prisma/seed.ts
    |   +-- src/ index, config, db, schemas, tokens, user-mapper,
    |            auth.service, users.service, routes, schemas.test
    +-- catalog-course-service/   prisma/schema.prisma + src/index.ts (routes 501)
    +-- exam-submission-service/  idem
    +-- analytics-streak-service/ idem + src/streak.ts (+ tests)
    +-- admin-contact-service/    idem
    +-- adaptive-qcm-service/     idem + src/adaptive-engine.ts (+ tests)
\end{lstlisting}

\begin{table}[H]\centering\small
\begin{tabularx}{\linewidth}{l X}
\toprule
\textbf{Commande (dans \code{backend/})} & \textbf{Effet}\\
\midrule
\code{npm run infra:up} & Démarre PostgreSQL, Redis et MinIO avec Docker Compose\\
\code{npm run db:generate} & Génère les 6 clients Prisma (\code{src/generated/prisma} de chaque service)\\
\code{npm run db:push} & Crée ou met à jour les 6 schémas PostgreSQL\\
\code{npm run db:seed} & Crée les comptes de démonstration (admin et étudiant)\\
\code{npm run dev} & Lance les 7 processus en parallèle (\code{concurrently} + \code{tsx watch})\\
\code{npm run typecheck} & \code{tsc} strict sur chaque workspace\\
\code{npm test} & Tests unitaires (\code{node:test} via \code{tsx})\\
\bottomrule
\end{tabularx}
\caption{Scripts du monorepo}
\end{table}

\begin{note}[Choix techniques]
\begin{itemize}
  \item \textbf{Express 5} plutôt que NestJS (le CDC accepte les deux). Express est plus simple à relire, et sa version 5 transmet automatiquement les erreurs des fonctions \code{async} au gestionnaire d'erreurs, sans \code{try/catch} dans chaque route.
  \item \textbf{tsx} exécute le TypeScript directement, sans étape de compilation en développement. \texttt{tsc -{}-noEmit} sert uniquement à la vérification de types.
  \item \textbf{Zod} valide toutes les entrées. Une donnée invalide produit une erreur \code{VALIDATION_ERROR} (400) qui détaille chaque champ fautif.
  \item \textbf{Un fichier \code{.env} unique} dans \code{backend/}, chargé par \code{@cnc/common} avec \code{process.loadEnvFile} (fonction native de Node). Chaque service y lit sa propre URL de base (\code{AUTH_DATABASE_URL}, \code{CATALOG_DATABASE_URL}, etc.).
\end{itemize}
\end{note}

% -----------------------------------------------------------------------------
\subsection{Le paquet partagé \texttt{@cnc/common}}

Tous les services sont construits de la même façon, grâce à \code{createServiceApp}. Pour relire un service, il suffit donc de lire son fichier \code{routes} : le reste est commun.

\begin{lstlisting}[language=TypeScript,caption={\texttt{packages/common/src/service-app.ts} (extrait)}]
export function createServiceApp({ name, internalSecret, router }: ServiceOptions): Express {
  const app = express()
  app.disable('x-powered-by')

  app.get('/health', (_req, res) => {                 // public : supervision
    res.json({ service: name, status: 'ok', time: new Date().toISOString() })
  })

  app.use(requireGateway(internalSecret))             // 403 si l'appel ne vient pas de la gateway
  app.use(express.json({ limit: '1mb' }))
  app.use(router)                                     // routes métier du service

  app.use(notFoundHandler)                            // 404 JSON
  app.use(errorHandler)                               // HttpError, ZodError, 500
  return app
}
\end{lstlisting}

\begin{lstlisting}[language=TypeScript,caption={\texttt{auth-context.ts} : contrôle d'accès dans les services}]
export function requireUser(role?: Role): RequestHandler {
  return (req, _res, next) => {
    const user = getAuthUser(req)                     // lit x-user-id / x-user-role
    if (!user) return next(new HttpError(401, 'Authentification requise', 'UNAUTHORIZED'))
    if (role && user.role !== role) return next(new HttpError(403, 'Droits insuffisants', 'FORBIDDEN'))
    next()
  }
}
\end{lstlisting}

Toutes les erreurs ont le même format JSON, quel que soit le service. Le frontend n'aura donc qu'un seul cas à gérer :
\begin{lstlisting}[numbers=none,language={}]
{ "error": "EMAIL_TAKEN", "message": "Un compte existe deja avec cet e-mail" }
{ "error": "VALIDATION_ERROR", "message": "Donnees invalides",
  "details": [ { "path": "filiere", "message": "Invalid option: ..." } ] }
\end{lstlisting}

% -----------------------------------------------------------------------------
\subsection{L'API Gateway}

\subsubsection{Chaîne de traitement d'une requête}

\begin{figure}[H]\centering
\begin{tikzpicture}[node distance=3.5mm,
  st/.style={box,fill=navy!8,minimum width=2.25cm,minimum height=1.25cm,font=\scriptsize,text width=2cm}]
  \node[st] (1) {\textbf{CORS}\\origine = \texttt{CORS\_ORIGIN}};
  \node[st,right=of 1] (2) {\textbf{Rate limit}\\600 req / 15\,min\\(20 sur login)};
  \node[st,right=of 2] (3) {\textbf{/api/health}\\santé agrégée des 6 services};
  \node[st,right=of 3,fill=warm!10,draw=warm] (4) {\textbf{authenticate}\\JWT + en-têtes d'identité};
  \node[st,right=of 4] (5) {\textbf{Proxy}\\préfixe $\rightarrow$ service\\\texttt{/api} retiré};
  \foreach \a/\b in {1/2,2/3,3/4,4/5}{\draw[arr] (\a) -- (\b);}
  \node[box,fill=red!6,draw=warm,font=\scriptsize,below=7mm of 4] (e) {401\\\texttt{UNAUTHORIZED}\\\texttt{INVALID\_TOKEN}};
  \node[box,fill=red!6,draw=warm,font=\scriptsize,below=7mm of 5] (e2) {503\\\texttt{SERVICE\_}\\\texttt{UNAVAILABLE}};
  \node[box,fill=red!6,draw=warm,font=\scriptsize,below=7mm of 2] (e3) {429\\Too Many\\Requests};
  \draw[darr,color=warm] (4) -- (e); \draw[darr,color=warm] (5) -- (e2); \draw[darr,color=warm] (2) -- (e3);
\end{tikzpicture}
\caption{Middlewares de la gateway, dans l'ordre d'exécution}
\end{figure}

La gateway \textbf{ne lit pas le corps} des requêtes (pas de \code{express.json()}). Le flux est transmis tel quel au service, ce qui évite de devoir le reconstruire dans le proxy.

\subsubsection{Table de routage}

\begin{table}[H]\centering\small
\begin{tabular}{l l l l}
\toprule
\textbf{Préfixe public} & \textbf{Service} & \textbf{Cible} & \textbf{Exemple réécrit}\\
\midrule
\code{/api/auth} & auth-user-service & \code{:4001} & \code{/api/auth/login} $\rightarrow$ \code{/auth/login}\\
\code{/api/users} & auth-user-service & \code{:4001} & \code{/api/users/me} $\rightarrow$ \code{/users/me}\\
\code{/api/catalog} & catalog-course-service & \code{:4002} & \code{/catalog/chapters}\\
\code{/api/exams} & exam-submission-service & \code{:4003} & \code{/exams/:id/sessions}\\
\code{/api/analytics} & analytics-streak-service & \code{:4004} & \code{/analytics/me/streak}\\
\code{/api/contact} & admin-contact-service & \code{:4005} & \code{/contact/tickets}\\
\code{/api/qcm} & adaptive-qcm-service & \code{:4006} & \code{/qcm/sessions}\\
\bottomrule
\end{tabular}
\caption{Routage de la gateway (\texttt{gateway/src/routes.ts})}
\end{table}

\textbf{Routes publiques} (sans jeton) : \code{POST /api/auth/register}, \code{/login}, \code{/refresh}, \code{/logout} et \code{GET /api/health}. Toutes les autres routes exigent un access token valide.

\begin{lstlisting}[language=TypeScript,caption={\texttt{gateway/src/authenticate.ts} : le cœur de la sécurité}]
export function authenticate(accessSecret: string, internalSecret: string): RequestHandler {
  return (req, _res, next) => {
    delete req.headers[HEADER_USER_ID]                // anti-usurpation : on ignore
    delete req.headers[HEADER_USER_ROLE]              // ce que le client prétend être
    req.headers[HEADER_INTERNAL_SECRET] = internalSecret

    const header = req.header('authorization')
    const token = header?.startsWith('Bearer ') ? header.slice(7) : null
    if (token) {
      try {
        const payload = jwt.verify(token, accessSecret) as AccessTokenPayload
        req.headers[HEADER_USER_ID] = payload.sub
        req.headers[HEADER_USER_ROLE] = payload.role
      } catch {
        return next(new HttpError(401, 'Access token invalide ou expire', 'INVALID_TOKEN'))
      }
    }
    if (!token && !isPublicRoute(req.method, req.path)) {
      return next(new HttpError(401, 'Authentification requise', 'UNAUTHORIZED'))
    }
    next()
  }
}
\end{lstlisting}

% -----------------------------------------------------------------------------
\subsection{auth-user-service (Sprint 1, implémenté)}\label{sec:auth}

\subsubsection{Découpage en couches}

\begin{figure}[H]\centering
\begin{tikzpicture}[node distance=4mm,
  ly/.style={box,minimum width=3.1cm,minimum height=1cm,font=\scriptsize,text width=2.9cm}]
  \node[ly,fill=royal!8] (r) {\textbf{routes.ts}\\HTTP : lit \texttt{req}, répond};
  \node[ly,fill=yellow!10,right=of r] (z) {\textbf{schemas.ts}\\validation Zod};
  \node[ly,fill=green!8,right=of z] (s) {\textbf{auth.service.ts}\\\textbf{users.service.ts}\\règles métier};
  \node[ly,fill=navy!8,right=of s] (d) {\textbf{db.ts}\\client Prisma};
  \node[db,right=6mm of d,minimum width=1.6cm] (pg) {schéma\\\texttt{auth}};
  \draw[arr] (r) -- (z); \draw[arr] (z) -- (s); \draw[arr] (s) -- (d); \draw[arr] (d) -- (pg);
  \node[ly,fill=soft,below=6mm of s] (t) {\textbf{tokens.ts}\\JWT, refresh, SHA-256};
  \node[ly,fill=soft,below=6mm of z] (m) {\textbf{user-mapper.ts}\\\texttt{toPublicUser()}};
  \node[ly,fill=soft,below=6mm of d] (c) {\textbf{config.ts}\\secrets, durées};
  \draw[arr] (s) -- (t); \draw[arr] (s) -- (m); \draw[darr] (t) -- (c);
\end{tikzpicture}
\caption{Couches du service : chaque fichier a une seule responsabilité}
\end{figure}

\subsubsection{Modèle de données (schéma \texttt{auth})}

Le modèle traduit la hiérarchie \code{User} / \code{Student} / \code{Administrator} du diagramme de classes du CDC (section 4.2). Comme Prisma ne gère pas l'héritage, on utilise une \textbf{table par spécialisation}, chacune en relation 1–1 avec \code{users}.

\begin{figure}[H]\centering
\begin{tikzpicture}[
  ent/.style={draw=navy,rectangle split,rectangle split parts=2,rectangle split part fill={navy!15,white},font=\scriptsize,align=left,rounded corners=2pt,text width=3.9cm}]
  \node[ent] (u) at (0,0) {\textbf{users}\nodepart{two}\texttt{id} UUID PK\\ \texttt{firstName}, \texttt{lastName}\\ \texttt{email} UNIQUE\\ \texttt{passwordHash} (bcrypt)\\ \texttt{role} STUDENT|ADMIN\\ \texttt{createdAt}, \texttt{updatedAt}};
  \node[ent] (s) at (-5.2,-0.4) {\textbf{students}\nodepart{two}\texttt{userId} PK/FK\\ \texttt{filiere} MP|PSI|TSI\\ \texttt{enrolledYear} SUP|SPE\\ \texttt{cpgeCenter}\\ \texttt{accessAllYears} = true};
  \node[ent] (a) at (5.2,-0.4) {\textbf{administrators}\nodepart{two}\texttt{userId} PK/FK\\ \texttt{privilegeLevel}\\ \texttt{assignedSubjects[]}\\ \texttt{department}};
  \node[ent] (t) at (0,-4.1) {\textbf{refresh\_tokens}\nodepart{two}\texttt{id} UUID PK\\ \texttt{userId} FK (index)\\ \texttt{tokenHash} UNIQUE (SHA-256)\\ \texttt{expiresAt}, \texttt{revokedAt?}};
  \draw[arr,-] (s) -- node[lbl,above]{1 — 0..1} (u);
  \draw[arr,-] (a) -- node[lbl,above]{0..1 — 1} (u);
  \draw[arr,-] (u) -- node[lbl,right]{1 — 0..*} (t);
\end{tikzpicture}
\caption{Schéma \texttt{auth} (suppression en cascade depuis \texttt{users})}
\end{figure}

\subsubsection{Endpoints}

\begin{table}[H]\centering\small
\begin{tabularx}{\linewidth}{l l l X}
\toprule
\textbf{Méthode} & \textbf{Route (gateway)} & \textbf{Accès} & \textbf{Corps / réponse}\\
\midrule
POST & \code{/api/auth/register} & public & \code{firstName, lastName, email, password, cpgeCenter, enrolledYear, filiere} $\rightarrow$ 201 \code{{user, accessToken, refreshToken}}\\
POST & \code{/api/auth/login} & public & \code{email, password} $\rightarrow$ 200 \code{{user, accessToken, refreshToken}}\\
POST & \code{/api/auth/refresh} & public & \code{refreshToken} $\rightarrow$ nouvelle paire (rotation)\\
POST & \code{/api/auth/logout} & public & \code{refreshToken} $\rightarrow$ 204 (jeton révoqué)\\
GET & \code{/api/users/me} & connecté & profil complet (jamais \code{passwordHash})\\
PATCH & \code{/api/users/me} & connecté & nom, prénom ; centre, année, filière (étudiant)\\
GET & \code{/api/users?page=&pageSize=} & ADMIN & liste paginée \code{{total, page, items}}\\
\bottomrule
\end{tabularx}
\caption{API de \texttt{auth-user-service}}
\end{table}

\subsubsection{Stratégie des jetons}

\begin{table}[H]\centering\small
\begin{tabularx}{\linewidth}{l X X}
\toprule
 & \textbf{Access token} & \textbf{Refresh token}\\
\midrule
Format & JWT signé HS256, contenu \code{{sub, role}} & 48 octets aléatoires (\code{base64url}), opaque\\
Durée & 15 minutes (\code{ACCESS_TOKEN_TTL}) & 30 jours (\code{REFRESH_TOKEN_TTL_DAYS})\\
Stockage serveur & aucun (vérifié par signature) & seulement son \textbf{empreinte SHA-256} en base\\
Vérifié par & la gateway, à chaque requête & \code{auth-user-service}, sur \code{/auth/refresh}\\
Révocation & expiration naturelle & \code{revokedAt} (déconnexion, rotation, vol détecté)\\
\bottomrule
\end{tabularx}
\caption{Les deux jetons définis par le CDC (section 2.1)}
\end{table}

\subsubsection{Rotation des refresh tokens et détection de vol}

Chaque appel à \code{/auth/refresh} \textbf{consomme} le refresh token présenté et en émet un nouveau. Si un jeton déjà consommé est présenté une deuxième fois, il a forcément été copié. Le service révoque alors \emph{tous} les jetons de l'utilisateur, ce qui coupe aussi la session de l'attaquant.

\begin{figure}[H]\centering
\begin{tikzpicture}[node distance=9mm]
  \node[box,fill=white] (in) {\texttt{POST /auth/refresh}\\\scriptsize\texttt{\{refreshToken\}}};
  \node[box,fill=soft,below=of in] (h) {recherche par \texttt{sha256(token)}};
  \node[diamond,draw=navy,aspect=2.2,font=\scriptsize,inner sep=1pt,below=of h] (q1) {trouvé ?};
  \node[diamond,draw=navy,aspect=2.2,font=\scriptsize,inner sep=1pt,below=of q1] (q2) {révoqué ?};
  \node[diamond,draw=navy,aspect=2.2,font=\scriptsize,inner sep=1pt,below=of q2] (q3) {expiré ?};
  \node[done,below=of q3] (ok) {révoque l'ancien + émet\\une nouvelle paire $\rightarrow$ 200};
  \node[box,fill=red!8,draw=warm,right=18mm of q1] (e1) {401 \texttt{INVALID\_REFRESH\_TOKEN}};
  \node[box,fill=red!15,draw=warm,right=18mm of q2,text width=4.3cm] (e2) {\textbf{vol présumé} : révoque \emph{tous} les jetons de l'utilisateur $\rightarrow$ 401 \texttt{REFRESH\_TOKEN\_REUSED}};
  \node[box,fill=red!8,draw=warm,right=18mm of q3] (e3) {401 \texttt{REFRESH\_TOKEN\_EXPIRED}};
  \draw[arr] (in) -- (h); \draw[arr] (h) -- (q1);
  \draw[arr] (q1) -- node[lbl,left]{oui} (q2); \draw[arr] (q2) -- node[lbl,left]{non} (q3); \draw[arr] (q3) -- node[lbl,left]{non} (ok);
  \draw[arr] (q1) -- node[lbl,above]{non} (e1); \draw[arr] (q2) -- node[lbl,above]{oui} (e2); \draw[arr] (q3) -- node[lbl,above]{oui} (e3);
\end{tikzpicture}
\caption{Algorithme de \texttt{refresh()} dans \texttt{auth.service.ts}}
\end{figure}

\begin{lstlisting}[language=TypeScript,caption={\texttt{auth.service.ts} : inscription et connexion}]
export async function register(input: RegisterInput) {
  const exists = await prisma.user.findUnique({ where: { email: input.email } })
  if (exists) throw new HttpError(409, 'Un compte existe deja avec cet e-mail', 'EMAIL_TAKEN')

  const user = await prisma.user.create({
    data: {
      firstName: input.firstName, lastName: input.lastName, email: input.email,
      passwordHash: await bcrypt.hash(input.password, config.bcryptRounds),   // coût 12
      role: 'STUDENT',
      student: { create: { filiere: input.filiere, enrolledYear: input.enrolledYear,
                           cpgeCenter: input.cpgeCenter, accessAllYears: true } },
    },
    include: withProfile,
  })
  return { user: toPublicUser(user), ...(await issueTokens(user.id, user.role)) }
}

export async function login(input: LoginInput) {
  const user = await prisma.user.findUnique({ where: { email: input.email }, include: withProfile })
  // compare() est appelé même si l'utilisateur n'existe pas (temps de réponse constant)
  const valid = await bcrypt.compare(input.password, user?.passwordHash ?? DUMMY_HASH)
  if (!user || !valid) throw new HttpError(401, 'E-mail ou mot de passe incorrect', 'INVALID_CREDENTIALS')
  return { user: toPublicUser(user), ...(await issueTokens(user.id, user.role)) }
}
\end{lstlisting}

\subsubsection{Mesures de sécurité}

\begin{table}[H]\centering\small
\begin{tabularx}{\linewidth}{l X}
\toprule
\textbf{Menace} & \textbf{Parade mise en place}\\
\midrule
Vol de la base de données & bcrypt (coût 12) pour les mots de passe, SHA-256 pour les refresh tokens\\
Brute force sur la connexion & 20 tentatives par tranche de 15 min et par IP sur login et register\\
Énumération des comptes & message d'erreur identique, \code{bcrypt.compare} exécuté même si l'e-mail n'existe pas\\
Usurpation d'identité & en-têtes \code{x-user-*} venant du client supprimés à la gateway\\
Appel direct d'un service & secret interne obligatoire, sinon 403\\
Vol d'un refresh token & rotation à chaque usage et révocation globale en cas de réutilisation\\
Fuite de données & \code{toPublicUser()} n'expose jamais \code{passwordHash} ; en-tête \code{x-powered-by} désactivé\\
Élévation de privilèges & le rôle vient du JWT signé ; \code{requireUser('ADMIN')} sur les routes d'administration\\
Entrées malformées & validation Zod de tous les corps de requête (e-mail normalisé, énumérations strictes)\\
\bottomrule
\end{tabularx}
\caption{Sécurité implémentée au Sprint 1}
\end{table}

% -----------------------------------------------------------------------------
\subsection{Les cinq autres microservices (squelettes)}

Chaque squelette contient trois éléments :
\begin{enumerate}
  \item un \textbf{schéma Prisma complet}, déjà poussé en base, qui traduit le diagramme de classes du CDC ;
  \item un fichier \code{index.ts} qui \textbf{déclare toutes les routes prévues}. Elles répondent \code{501 NOT_IMPLEMENTED} avec le numéro du sprint concerné ;
  \item le contrôle d'accès déjà en place sur les routes d'administration (\code{requireUser('ADMIN')}).
\end{enumerate}

\begin{figure}[H]\centering
\begin{tikzpicture}[
  ent/.style={draw=navy,rectangle split,rectangle split parts=2,rectangle split part fill={#1!15,white},font=\tiny,align=left,rounded corners=2pt,text width=3.55cm},
  ent/.default=royal]
  \node[font=\small\bfseries\color{royal}] at (0,0.6) {catalog};
  \node[ent,anchor=north] (sb) at (0,0.3) {\textbf{Subject}\nodepart{two}code MATH|PHYS|SI|INFO\\coeffMP, coeffPSI, coeffTSI};
  \node[ent,below=2mm of sb] (ch) {\textbf{Chapter}\nodepart{two}studyLevel SUP|SPE, domainName\\title, targetFilieres[], orderIndex};
  \node[ent,below=2mm of ch] (ls) {\textbf{Lesson}\nodepart{two}contentLatex, coursePdfUrl\\summarySheetUrl, durée};
  \node[ent,below=2mm of ls] (ex) {\textbf{Exercise}\nodepart{two}ref (EXO-MP-MATH-204), difficulty\\statementLatex, hintsLatex[]\\solutionLatex, statementPdfUrl, corrPdfUrl};
  \node[ent,below=2mm of ex] (ea) {\textbf{ExerciseAttempt}\nodepart{two}usedMode PDF|BIBMATH, hintUsed\\resolutionStatus, timeSpentSec};
  \node[ent,below=2mm of ea] (rb) {\textbf{ReferenceBook}\nodepart{two}title, authors, edition\\studyLevel SUP|SPE|BOTH};

  \node[font=\small\bfseries\color{warm}] at (4.1,0.6) {exam};
  \node[ent=warm,anchor=north] (ce) at (4.1,0.3) {\textbf{CncExam}\nodepart{two}kind CNC\_ANNALE|PROBLEM\\year, filiere, epreuve\\durationMinutes, PDFs, barème, rapport jury};
  \node[ent=warm,below=2mm of ce] (es) {\textbf{ExamSession}\nodepart{two}startedAt, deadlineAt\\submittedAt (chronomètre)};
  \node[ent=warm,below=2mm of es] (ws) {\textbf{WrittenSubmission}\nodepart{two}writtenContentLatex\\uploadedScanUrl, selfGrade\\personalNotes};

  \node[font=\small\bfseries\color{okgreen}] at (8.2,0.6) {analytics};
  \node[ent=okgreen,anchor=north] (us) at (8.2,0.3) {\textbf{UserStreak}\nodepart{two}currentStreak, longestStreak\\lastActiveDate, freezeTokens};
  \node[ent=okgreen,below=2mm of us] (dl) {\textbf{DailyActivityLog}\nodepart{two}(userId, date, subject) UNIQUE\\lessonsRead, exercisesSolved\\problemsWritten, qcmCompleted, minutes};
  \node[font=\small\bfseries\color{navy}] at (8.2,-3.4) {contact};
  \node[ent=navy,anchor=north] (cm) at (8.2,-3.7) {\textbf{ContactMessage}\nodepart{two}category (4 valeurs), targetEntityRef\\status OPEN|IN\_PROGRESS|RESOLVED\\adminReply, repliedBy};
  \node[ent=navy,below=2mm of cm] (nt) {\textbf{Notification}\nodepart{two}ticketId, message, readAt};

  \node[font=\small\bfseries\color{purple}] at (12.3,0.6) {qcm};
  \node[ent=purple,anchor=north] (qq) at (12.3,0.3) {\textbf{QcmQuestion}\nodepart{two}questionLatex, difficultyLevel\\targetTimeSec, explanationLatex};
  \node[ent=purple,below=2mm of qq] (qc) {\textbf{QcmChoice}\nodepart{two}labelLatex, isCorrect, feedback};
  \node[ent=purple,below=2mm of qc] (qm) {\textbf{QcmMastery}\nodepart{two}(userId, chapterId), score = 50};
  \node[ent=purple,below=2mm of qm] (qs) {\textbf{QcmSession} / \textbf{QcmAnswer}\nodepart{two}currentDifficulty, score\\weakLessonsJson, masteryAfter};
  \node[ent=purple,below=2mm of qs] (qr) {\textbf{QcmSpacedRepetition}\nodepart{two}step, intervalDays\\easeFactor, nextReviewDate};
\end{tikzpicture}
\caption{Modèles Prisma des services squelettes (déjà créés en base)}
\end{figure}

\begin{table}[H]\centering\footnotesize
\begin{tabularx}{\linewidth}{l X}
\toprule
\textbf{Service} & \textbf{Routes déclarées (réponse 501 en attendant)}\\
\midrule
catalog (S2) & \code{GET /catalog/subjects}, \code{/chapters?subject&level&filiere}, \code{/chapters/:id}, \code{/lessons/:id}, \code{/lessons/:id/exercises}, \code{/exercises/:id}, \code{POST /exercises/:id/attempts}, \code{GET /books}, \code{/books/:id} ; admin : \code{POST|PATCH|DELETE /catalog/:entity}\\
exam (S3) & \code{GET /exams?kind&year&filiere&epreuve}, \code{/exams/:id}, \code{POST /exams/:id/sessions}, \code{POST /exams/sessions/:id/submit}, \code{PATCH /exams/submissions/:id}, \code{GET /exams/submissions/mine} ; admin : \code{POST|PATCH /exams}\\
analytics (S4) & \code{POST /analytics/activities}, \code{GET /analytics/me/streak}, \code{/me/heatmap?year}, \code{/me/subjects}, \code{/me/summary}\\
contact (S4) & \code{POST /contact/tickets}, \code{GET /tickets/mine}, \code{/tickets/:id}, \code{/notifications}, \code{PATCH /notifications/:id/read} ; admin : \code{GET /contact/admin/tickets}, \code{PATCH /admin/tickets/:id}, \code{POST /admin/tickets/:id/reply}\\
qcm (S5) & \code{POST /qcm/sessions}, \code{POST /sessions/:id/answers}, \code{POST /sessions/:id/finish}, \code{GET /qcm/flash}, \code{GET /qcm/mastery} ; admin : \code{POST|PATCH /qcm/questions}\\
\bottomrule
\end{tabularx}
\caption{Contrat d'API déjà fixé pour les sprints 2 à 5}
\end{table}

% -----------------------------------------------------------------------------
\subsection{Logique métier déjà codée : Streak et QCM adaptatif}

Les deux algorithmes du CDC sont écrits comme des \textbf{fonctions pures}, sans base de données. Ils sont donc testables en isolation ; les sprints 4 et 5 n'auront qu'à les appeler depuis les routes.

\subsubsection{Moteur de Streak (\texttt{analytics-streak-service/src/streak.ts})}

Le calcul suit le diagramme d'activité 4.3.1 du CDC. On note $\Delta$ le nombre de jours calendaires entre la dernière activité et aujourd'hui :
\[
\text{currentStreak}' =
\begin{cases}
\text{currentStreak} & \text{si } \Delta = 0 \quad\text{(déjà actif aujourd'hui)}\\
\text{currentStreak}+1 & \text{si } \Delta = 1\\
1 & \text{si } \Delta > 1
\end{cases}
\]
\[
\text{longestStreak}' = \max(\text{longestStreak},\ \text{currentStreak}')
\]

\begin{lstlisting}[language=TypeScript,caption={Fonction \texttt{updateStreak}}]
export function updateStreak(state: StreakState, today: string): StreakState {
  if (state.lastActiveDate === null) {
    return { currentStreak: 1, longestStreak: Math.max(1, state.longestStreak), lastActiveDate: today }
  }
  const delta = daysBetween(state.lastActiveDate, today)
  if (delta < 0) throw new Error('Activite anterieure a la derniere activite enregistree')
  if (delta === 0) return state
  const currentStreak = delta === 1 ? state.currentStreak + 1 : 1
  return { currentStreak, longestStreak: Math.max(state.longestStreak, currentStreak), lastActiveDate: today }
}
\end{lstlisting}

\begin{note}[Le piège du fuseau horaire]
Un « jour » doit être le jour \emph{de l'étudiant} (Africa/Casablanca, UTC+1), pas celui du serveur. \code{localDate()} convertit un instant en date locale : une activité à 23h30 UTC le 4 novembre compte pour le \textbf{5} novembre à Casablanca. Un test unitaire couvre ce cas.
\end{note}

\subsubsection{Moteur QCM adaptatif (\texttt{adaptive-qcm-service/src/adaptive-engine.ts})}

Formule (1) du CDC (section 6.2) :
\[
M_{k+1}(u,c) = \operatorname{clamp}_{[0,100]}\Big(M_k(u,c) + \alpha\cdot S_k\cdot\omega(D_k)\cdot\varphi(t_k)\Big)
\]
\[
S_k=\begin{cases}+1 & \text{réponse exacte}\\-1{,}2 & \text{réponse fausse}\end{cases}
\qquad
\omega(D)=1+0{,}25\,(D-1)
\qquad
\varphi(t)=\max\!\Big(0{,}5\,;\ \min\big(1{,}3\,;\ \tfrac{T_{\text{cible}}}{t}\big)\Big)
\]
Le CDC ne fixe pas la valeur de $\alpha$. Elle est réglée à $\alpha=5$ (constante \code{ALPHA}) : une bonne réponse rapide en difficulté 4 rapporte au plus $5\times1{,}75\times1{,}3\approx 11{,}4$ points.

\begin{figure}[H]\centering
\begin{tikzpicture}[node distance=7mm]
  \node[box] (q) {question de niveau $D$};
  \node[diamond,draw=navy,aspect=2.4,font=\scriptsize,right=of q,inner sep=1pt] (c) {exacte ?};
  \node[diamond,draw=navy,aspect=2.4,font=\scriptsize,above right=2mm and 8mm of c,inner sep=1pt] (f) {$t\le T_{\text{cible}}$ ?};
  \node[done,right=8mm of f,yshift=5mm] (up) {$D\leftarrow\min(4,D+1)$};
  \node[box,right=8mm of f,yshift=-5mm] (eq) {$D$ inchangé};
  \node[box,fill=red!8,draw=warm,below right=2mm and 8mm of c,text width=4.2cm,font=\scriptsize] (dn) {$D\leftarrow\max(1,D-1)$, explication affichée, ajout en répétition espacée (J+1)};
  \draw[arr] (q) -- (c);
  \draw[arr] (c) |- node[lbl,pos=0.25,left]{oui} (f);
  \draw[arr] (c) |- node[lbl,pos=0.25,left]{non} (dn);
  \draw[arr] (f) -- node[lbl,above]{oui} (up);
  \draw[arr] (f) -- node[lbl,below]{non} (eq);
\end{tikzpicture}
\caption{Sélection adaptative de la difficulté (\texttt{nextDifficulty})}
\end{figure}

\paragraph{Répétition espacée.} Une question ratée entre dans la file avec une révision à J+1. Chaque bonne réponse lors d'une révision la fait avancer d'un cran ; une erreur la ramène à J+1.
\begin{center}
\begin{tikzpicture}[x=1cm]
  \draw[arr] (0,0) -- (12.6,0) node[right,font=\scriptsize]{temps};
  \foreach \x/\l/\s in {0/{erreur}/{J0},2.2/{révision 1}/{J+1},4.8/{révision 2}/{J+3},7.8/{révision 3}/{J+7}}{
    \fill[navy] (\x,0) circle (2.5pt);
    \node[font=\scriptsize,above] at (\x,0.12) {\l};
    \node[font=\scriptsize\bfseries,below] at (\x,-0.12) {\s};}
  \fill[okgreen] (11.2,0) circle (3pt);
  \node[font=\scriptsize,above,text=okgreen] at (11.2,0.12) {révision 4 réussie};
  \node[font=\scriptsize\bfseries,below,text=okgreen] at (11.2,-0.12) {J+21 $\rightarrow$ sortie de la file};
\end{tikzpicture}
\end{center}

% -----------------------------------------------------------------------------
\subsection{Vérifications réalisées}

\begin{ok}[Compilation et tests unitaires]
\begin{itemize}
  \item \code{npm run typecheck} : \textbf{0 erreur} sur les 8 workspaces (TypeScript \code{strict}).
  \item \code{npm test} : \textbf{15 tests réussis sur 15}. Ils couvrent le Streak (6 tests : $\Delta=0$, $\Delta=1$, $\Delta>1$, changement de mois ou d'année, année bissextile, fuseau horaire), le moteur QCM (6 tests : bornes de $M$, de $\varphi$ et de $D$, séquence J+1/J+3/J+7/J+21) et la validation de l'inscription (3 tests).
  \item Frontend : \code{next build} réussi après le déplacement (11 pages générées).
\end{itemize}
\end{ok}

\paragraph{Test de bout en bout.} PostgreSQL a été lancé dans Docker, les 6 schémas créés, puis les 7 processus démarrés. Le scénario ci-dessous a été joué contre la gateway (\code{http://localhost:4000}) :

\begin{longtable}{>{\ttfamily}c p{6.3cm} p{6.6cm}}
\toprule
\textrm{\textbf{Code}} & \textbf{Requête} & \textbf{Comportement vérifié}\\
\midrule\endhead
201 & \code{POST /api/auth/register} & compte créé, rôle STUDENT, filière PSI, \code{passwordHash} absent de la réponse\\
409 & même e-mail une seconde fois & \code{EMAIL_TAKEN}\\
400 & e-mail invalide, filière « BCPST » & \code{VALIDATION_ERROR} avec 7 champs détaillés\\
401 & connexion avec un mauvais mot de passe & \code{INVALID_CREDENTIALS}\\
200 & connexion correcte & paire de jetons\\
401 & \code{GET /api/users/me} sans jeton & bloqué par la gateway\\
401 & \code{GET /api/users/me} avec un jeton forgé & \code{INVALID_TOKEN}\\
200 & \code{GET /api/users/me} avec un jeton valide & profil renvoyé\\
200 & \code{PATCH /api/users/me} \code{{enrolledYear: SPE}} & profil mis à jour\\
403 & \code{GET /api/users} en tant qu'étudiant & réservé à ADMIN\\
403 & même requête + en-tête \code{x-user-role: ADMIN} & \textbf{usurpation bloquée}\\
200 & \code{GET /api/users} en tant qu'admin & liste paginée (3 comptes)\\
200 & \code{POST /api/auth/refresh} & rotation : nouvelle paire\\
401 & ré-utilisation de l'ancien refresh token & \code{REFRESH_TOKEN_REUSED}\\
401 & refresh avec le jeton issu de la rotation & révoqué aussi (révocation globale)\\
204 & \code{POST /api/auth/logout} & jeton révoqué\\
501 & \code{GET /api/catalog/subjects} & « sera implémenté au Sprint 2 »\\
403 & appel direct \code{http://localhost:4001/users/me} & refusé sans secret interne\\
200 & \code{GET /api/health} & gateway + 6 services « up »\\
\bottomrule
\caption{Scénario de bout en bout : 19 cas, tous conformes}
\end{longtable}

\section{Base de données : état actuel et consignes}\label{sec:bdd}

Cette section décrit la base \textbf{telle qu'elle existe aujourd'hui}. Les chiffres ont été relevés directement dans PostgreSQL, et non recopiés depuis les schémas Prisma. Elle fixe ensuite les \textbf{règles à suivre} pour toute modification future.

% -----------------------------------------------------------------------------
\subsection{Vue d'ensemble}

\begin{table}[H]\centering\small
\begin{tabularx}{\linewidth}{>{\bfseries}l X}
\toprule
Moteur & PostgreSQL \textbf{16.15} (image Docker \code{postgres:16})\\
Conteneur & \code{backend-postgres-1}, volume persistant \code{nexora-postgres}\\
Connexion & \code{127.0.0.1:5434}, base \code{nexora}, utilisateur \code{nexora} (local uniquement)\\
Taille & 8,5\,Mo\\
Organisation & \textbf{6 schémas} logiques, un par microservice : \code{auth}, \code{catalog}, \code{exam}, \code{analytics}, \code{contact}, \code{qcm}\\
Objets & \textbf{23 tables}, \textbf{14 types énumérés}, \textbf{15 clés étrangères} (toutes internes à un schéma), \textbf{18 index} secondaires ou d'unicité\\
Outil & Prisma 6.19, un \code{schema.prisma} et un client généré par service\\
Méthode de création & \code{prisma db push} : il n'y a \textbf{pas encore d'historique de migrations} (aucune table \code{_prisma_migrations})\\
Cohérence & \code{prisma migrate diff} renvoie « No difference detected » pour les 6 services : la base correspond exactement aux fichiers \code{schema.prisma}\\
\bottomrule
\end{tabularx}
\caption{Fiche d'identité de la base de données actuelle}
\end{table}

\begin{figure}[H]\centering
\resizebox{\linewidth}{!}{%
\begin{tikzpicture}[
  sch/.style={draw=#1,fill=#1!6,rounded corners=4pt,inner sep=6pt,line width=0.8pt},
  tb/.style={draw=black!35,fill=white,font=\scriptsize\ttfamily,text width=3.1cm,align=left,inner sep=2.5pt,rounded corners=1pt},
  hd/.style={font=\small\bfseries},
  node distance=1.2mm]
  % auth
  \node[tb] (a1) at (0,0) {users \hfill\textrm{3}};
  \node[tb,below=of a1] (a2) {students \hfill\textrm{2}};
  \node[tb,below=of a2] (a3) {administrators \hfill\textrm{1}};
  \node[tb,below=of a3] (a4) {refresh\_tokens \hfill\textrm{4}};
  \node[hd,above=1.5mm of a1,text=okgreen] (ha) {auth \textmd{\scriptsize(Sprint 1)}};
  \begin{scope}[on background layer]\node[sch=okgreen,fit=(ha)(a1)(a4)] (A) {};\end{scope}
  % catalog
  \node[tb] (c1) at (4.2,0) {subjects};
  \node[tb,below=of c1] (c2) {chapters};
  \node[tb,below=of c2] (c3) {lessons};
  \node[tb,below=of c3] (c4) {exercises};
  \node[tb,below=of c4] (c5) {exercise\_attempts};
  \node[tb,below=of c5] (c6) {reference\_books};
  \node[hd,above=1.5mm of c1,text=royal] (hc) {catalog};
  \begin{scope}[on background layer]\node[sch=royal,fit=(hc)(c1)(c6)] (C) {};\end{scope}
  % exam
  \node[tb] (e1) at (8.4,0) {cnc\_exams};
  \node[tb,below=of e1] (e2) {exam\_sessions};
  \node[tb,below=of e2] (e3) {written\_submissions};
  \node[hd,above=1.5mm of e1,text=warm] (he) {exam};
  \begin{scope}[on background layer]\node[sch=warm,fit=(he)(e1)(e3)] (E) {};\end{scope}
  % analytics
  \node[tb] (n1) at (8.4,-2.9) {user\_streaks};
  \node[tb,below=of n1] (n2) {daily\_activity\_logs};
  \node[hd,above=1.5mm of n1,text=teal] (hn) {analytics};
  \begin{scope}[on background layer]\node[sch=teal,fit=(hn)(n1)(n2)] (N) {};\end{scope}
  % contact
  \node[tb] (t1) at (12.6,0) {contact\_messages};
  \node[tb,below=of t1] (t2) {notifications};
  \node[hd,above=1.5mm of t1,text=navy] (ht) {contact};
  \begin{scope}[on background layer]\node[sch=navy,fit=(ht)(t1)(t2)] (T) {};\end{scope}
  % qcm
  \node[tb] (q1) at (12.6,-2.2) {qcm\_questions};
  \node[tb,below=of q1] (q2) {qcm\_choices};
  \node[tb,below=of q2] (q3) {qcm\_mastery};
  \node[tb,below=of q3] (q4) {qcm\_sessions};
  \node[tb,below=of q4] (q5) {qcm\_answers};
  \node[tb,below=of q5] (q6) {qcm\_spaced\_repetition};
  \node[hd,above=1.5mm of q1,text=purple] (hq) {qcm};
  \begin{scope}[on background layer]\node[sch=purple,fit=(hq)(q1)(q6)] (Q) {};\end{scope}
  % references logiques userId
  \node[font=\scriptsize,align=left,anchor=north west] at (-1.7,-5.6) {%
    Les chiffres indiquent le nombre de lignes actuel (tables vides sinon).\\
    Références \textbf{logiques} par UUID, \textbf{sans} clé étrangère : toute colonne \texttt{userId} désigne \texttt{auth.users},\\ et \texttt{qcm} désigne \texttt{catalog} par \texttt{chapterId} et \texttt{lessonId}.};
\end{tikzpicture}}
\caption{Les 6 schémas et leurs 23 tables}
\end{figure}

% -----------------------------------------------------------------------------
\subsection{Contenu actuel}

Seul le schéma \code{auth} contient des données ; les cinq autres schémas ont leurs tables créées, mais vides.

\begin{table}[H]\centering\small
\begin{tabularx}{\linewidth}{l r X}
\toprule
\textbf{Table} & \textbf{Lignes} & \textbf{Origine}\\
\midrule
\code{auth.users} & 3 & 2 comptes de démonstration (\code{npm run db:seed}) + 1 compte \code{test.<horodatage>@cpge.ma} créé par le test de bout en bout\\
\code{auth.students} & 2 & Achraf Lemrani (MP, SPE) et le compte de test (PSI, SPE)\\
\code{auth.administrators} & 1 & Karim Mansouri (\code{privilegeLevel} 3, 4 matières)\\
\code{auth.refresh_tokens} & 4 & sessions du test ; \textbf{1 seule active}, les 3 autres révoquées (rotation, réutilisation, déconnexion)\\
autres schémas & 0 & tables créées, aucune donnée\\
\bottomrule
\end{tabularx}
\caption{Données présentes dans la base}
\end{table}

\begin{alerte}[À nettoyer]
Le compte créé par le test (\code{test.*@cpge.ma}) n'a pas d'intérêt : on peut le supprimer. La suppression en cascade emporte aussi son profil étudiant et ses jetons.
\begin{lstlisting}[numbers=none,language=SQL]
DELETE FROM auth.users WHERE email LIKE 'test.%@cpge.ma';
\end{lstlisting}
\end{alerte}

% -----------------------------------------------------------------------------
\subsection{Dictionnaire des données}

Notation : \textbf{PK} clé primaire, \textbf{FK} clé étrangère, \textbf{U} unique, \textbf{I} index, \textbf{?} colonne pouvant être nulle, \textbf{=} valeur par défaut.

{\footnotesize
\begin{longtable}{p{3cm} p{8.6cm} p{3.6cm}}
\toprule
\textbf{Table} & \textbf{Colonnes (type)} & \textbf{Clés / index}\\
\midrule\endhead
\multicolumn{3}{l}{\cellcolor{okgreen!12}\textbf{Schéma \texttt{auth}} — auth-user-service}\\
\code{users} & \code{id} uuid, \code{firstName} text, \code{lastName} text, \code{email} text, \code{passwordHash} text, \code{role} Role = STUDENT, \code{createdAt} timestamp = now, \code{updatedAt} timestamp & PK id ; U email\\
\code{students} & \code{userId} uuid, \code{filiere} Filiere, \code{enrolledYear} StudyYear, \code{cpgeCenter} text, \code{accessAllYears} bool = true & PK userId ; FK users (cascade)\\
\code{administrators} & \code{userId} uuid, \code{privilegeLevel} int = 1, \code{assignedSubjects} text[], \code{department} text? & PK userId ; FK users (cascade)\\
\code{refresh_tokens} & \code{id} uuid, \code{userId} uuid, \code{tokenHash} text, \code{expiresAt} timestamp, \code{revokedAt} timestamp?, \code{createdAt} timestamp = now & PK id ; U tokenHash ; I userId ; FK users (cascade)\\
\midrule
\multicolumn{3}{l}{\cellcolor{royal!12}\textbf{Schéma \texttt{catalog}} — catalog-course-service}\\
\code{subjects} & \code{id} int (auto-incrément), \code{code} SubjectCode, \code{name} text, \code{coeffMP}, \code{coeffPSI}, \code{coeffTSI} double & PK id ; U code\\
\code{chapters} & \code{id} uuid, \code{subjectId} int, \code{studyLevel} StudyLevel, \code{domainName} text, \code{title} text, \code{targetFilieres} text[], \code{orderIndex} int & PK ; I (subjectId, studyLevel) ; FK subjects (restrict)\\
\code{lessons} & \code{id} uuid, \code{chapterId} uuid, \code{title}, \code{contentLatex} text, \code{coursePdfUrl}?, \code{summarySheetUrl}? text, \code{estimatedDurationMin} int, \code{orderIndex} int & PK ; FK chapters (cascade)\\
\code{exercises} & \code{id} uuid, \code{ref} text, \code{lessonId} uuid, \code{annaleSource} text?, \code{difficulty} int, \code{statementLatex}, \code{solutionLatex} text, \code{hintsLatex} text[], \code{statementPdfUrl}?, \code{corrPdfUrl}? text & PK ; U ref ; FK lessons (cascade)\\
\code{exercise_attempts} & \code{id} uuid, \code{exerciseId} uuid, \code{userId} uuid, \code{usedMode} ExerciseMode, \code{hintUsed} bool = false, \code{resolutionStatus} ResolutionStatus, \code{timeSpentSec} int, \code{attemptedAt} timestamp = now & PK ; I userId ; FK exercises (cascade)\\
\code{reference_books} & \code{id} uuid, \code{subjectId} int, \code{title}, \code{authors} text, \code{edition} text?, \code{targetFilieres} text[], \code{studyLevel} BookLevel, \code{summaryAndAdvice} text, \code{resourcePdfUrl} text? & PK ; FK subjects (restrict)\\
\midrule
\multicolumn{3}{l}{\cellcolor{warm!12}\textbf{Schéma \texttt{exam}} — exam-submission-service}\\
\code{cnc_exams} & \code{id} uuid, \code{kind} ExamKind, \code{year} int?, \code{filiere} Filiere, \code{epreuve} Epreuve, \code{title} text, \code{durationMinutes} int, \code{statementPdfUrl} text, \code{solutionPdfUrl}?, \code{gradingSchemeUrl}?, \code{juryReportUrl}? text & PK ; I (year, filiere, epreuve)\\
\code{exam_sessions} & \code{id} uuid, \code{userId} uuid, \code{examId} uuid, \code{timed} bool = true, \code{startedAt} timestamp = now, \code{deadlineAt}?, \code{submittedAt}? timestamp & PK ; I userId ; FK cnc\_exams (restrict)\\
\code{written_submissions} & \code{id} uuid, \code{userId} uuid, \code{examId} uuid, \code{sessionId} uuid?, \code{writtenContentLatex}?, \code{uploadedScanUrl}? text, \code{selfGrade} double?, \code{personalNotes} text?, \code{submittedAt} timestamp = now & PK ; U sessionId ; I userId ; FK cnc\_exams (restrict), exam\_sessions (set null)\\
\midrule
\multicolumn{3}{l}{\cellcolor{teal!12}\textbf{Schéma \texttt{analytics}} — analytics-streak-service}\\
\code{user_streaks} & \code{userId} uuid, \code{currentStreak}, \code{longestStreak}, \code{freezeTokens} int = 0, \code{lastActiveDate} date? & PK userId\\
\code{daily_activity_logs} & \code{id} uuid, \code{userId} uuid, \code{activityDate} date, \code{subjectCode} SubjectCode, \code{lessonsRead}, \code{exercisesSolved}, \code{problemsWritten}, \code{qcmCompleted}, \code{studyMinutes} int = 0 & PK ; U (userId, activityDate, subjectCode)\\
\midrule
\multicolumn{3}{l}{\cellcolor{navy!12}\textbf{Schéma \texttt{contact}} — admin-contact-service}\\
\code{contact_messages} & \code{id} uuid, \code{userId} uuid, \code{category} TicketCategory, \code{targetEntityRef} text?, \code{messageBody} text, \code{status} TicketStatus = OPEN, \code{adminReply} text?, \code{repliedBy} uuid?, \code{createdAt} = now, \code{updatedAt}, \code{resolvedAt}? timestamp & PK ; I userId ; I status\\
\code{notifications} & \code{id} uuid, \code{userId} uuid, \code{ticketId} uuid, \code{message} text, \code{readAt} timestamp?, \code{createdAt} = now & PK ; I (userId, readAt) ; FK contact\_messages (cascade)\\
\midrule
\multicolumn{3}{l}{\cellcolor{purple!12}\textbf{Schéma \texttt{qcm}} — adaptive-qcm-service}\\
\code{qcm_questions} & \code{id} uuid, \code{lessonId} uuid, \code{chapterId} uuid, \code{questionLatex} text, \code{difficultyLevel} int, \code{targetTimeSec} int, \code{explanationLatex} text & PK ; I (chapterId, difficultyLevel)\\
\code{qcm_choices} & \code{id} uuid, \code{questionId} uuid, \code{labelLatex} text, \code{isCorrect} bool, \code{feedback} text? & PK ; FK qcm\_questions (cascade)\\
\code{qcm_mastery} & \code{userId} uuid, \code{chapterId} uuid, \code{masteryScore} double = 50, \code{updatedAt} timestamp & PK (userId, chapterId)\\
\code{qcm_sessions} & \code{id} uuid, \code{userId} uuid, \code{chapterId} uuid, \code{currentDifficulty} int = 2, \code{score} double?, \code{weakLessonsJson} jsonb?, \code{startedAt} = now, \code{endedAt}? timestamp & PK ; I userId\\
\code{qcm_answers} & \code{id} uuid, \code{sessionId}, \code{questionId}, \code{choiceId} uuid, \code{isCorrect} bool, \code{timeSpentSec} int, \code{masteryAfter} double, \code{answeredAt} = now & PK ; FK qcm\_sessions (cascade), qcm\_questions (restrict)\\
\code{qcm_spaced_}\allowbreak\code{repetition} & \code{userId}, \code{questionId} uuid, \code{step} int = 0, \code{intervalDays} int, \code{easeFactor} double = 2.5, \code{nextReviewDate} date & PK (userId, questionId) ; I (userId, nextReviewDate)\\
\bottomrule
\caption{Dictionnaire des données relevé dans PostgreSQL (23 tables)}
\end{longtable}
}

\begin{table}[H]\centering\footnotesize
\begin{tabularx}{\linewidth}{l l X}
\toprule
\textbf{Schéma} & \textbf{Type énuméré} & \textbf{Valeurs}\\
\midrule
auth & \code{Role} & STUDENT, ADMIN\\
auth & \code{Filiere} & MP, PSI, TSI\\
auth & \code{StudyYear} & SUP, SPE\\
catalog & \code{SubjectCode} & MATH, PHYS, SI, INFO\\
catalog & \code{StudyLevel} & SUP, SPE\\
catalog & \code{BookLevel} & SUP, SPE, BOTH\\
catalog & \code{ExerciseMode} & PDF, BIBMATH\\
catalog & \code{ResolutionStatus} & MASTERED, WITH\_HINT, TO\_REVIEW\\
exam & \code{ExamKind} & CNC\_ANNALE, PROBLEM\\
exam & \code{Filiere} & MP, PSI, TSI \hfill\emph{(copie de auth.Filiere)}\\
exam & \code{Epreuve} & MATH\_1, MATH\_2, PHYS\_1, PHYS\_2, SI, INFO\\
analytics & \code{SubjectCode} & MATH, PHYS, SI, INFO \hfill\emph{(copie de catalog.SubjectCode)}\\
contact & \code{TicketCategory} & ERRATUM, PEDAGOGICAL\_QUESTION, TECHNICAL\_ISSUE, BOOK\_SUGGESTION\\
contact & \code{TicketStatus} & OPEN, IN\_PROGRESS, RESOLVED\\
\bottomrule
\end{tabularx}
\caption{Les 14 types énumérés PostgreSQL}
\end{table}

\paragraph{Choix des règles de suppression.} \textbf{CASCADE} quand l'enfant n'a pas de sens sans son parent : profil et jetons d'un utilisateur, leçons d'un chapitre, choix d'une question. \textbf{RESTRICT} quand la suppression détruirait un historique : on ne peut pas supprimer une matière qui a des chapitres, ni une épreuve qui a des copies ou des sessions. \textbf{SET NULL} quand le lien est facultatif : une copie survit à la suppression de sa session chronométrée.

% -----------------------------------------------------------------------------
\subsection{Conventions observées dans la base}

\begin{itemize}
  \item \textbf{Noms de tables} en \code{snake_case} au pluriel (attribut \code{@@map} de Prisma). \textbf{Noms de colonnes} en \code{camelCase}, comportement par défaut de Prisma. En SQL écrit à la main, il faut donc des guillemets :
\begin{lstlisting}[numbers=none,language=SQL]
SELECT "firstName", "lastName", role FROM auth.users WHERE "createdAt" > now() - interval '7 days';
\end{lstlisting}
  \item \textbf{Les UUID sont générés par Prisma}, côté application. La colonne \code{id} n'a \emph{pas} de valeur par défaut en base, et \code{updatedAt} non plus. Une insertion SQL manuelle doit donc fournir ces valeurs, par exemple \code{gen_random_uuid()} et \code{now()}.
  \item \textbf{Dates et heures} : les colonnes \code{...At} sont des \code{timestamp without time zone} dans lesquels Prisma écrit en \textbf{UTC}. Les colonnes \code{...Date} (\code{activityDate}, \code{lastActiveDate}, \code{nextReviewDate}) sont de type \code{date} et représentent le \textbf{jour local de l'étudiant} (Africa/Casablanca).
  \item \textbf{Tableaux} (\code{text[]}) : ils peuvent être nuls en base, mais Prisma les traite comme obligatoires et les initialise à \texttt{\{\}}.
  \item \textbf{Énumérations dupliquées} : \code{Filiere} existe dans \code{auth} et dans \code{exam}, \code{SubjectCode} dans \code{catalog} et dans \code{analytics}. Cette duplication est voulue, pour que chaque service reste indépendant, mais les copies doivent être \textbf{maintenues identiques}.
  \item \textbf{Références entre services} (\code{userId}, \code{chapterId}, \code{lessonId}, \code{repliedBy}) : ce sont de simples UUID, sans clé étrangère. C'est l'application qui garantit leur validité.
\end{itemize}

% -----------------------------------------------------------------------------
\subsection{Limites actuelles}

\begin{enumerate}
  \item \textbf{Pas de migrations versionnées} : la base a été créée avec \code{db push}. On ne peut ni rejouer l'historique sur une autre machine, ni revenir en arrière proprement, et un \code{db push} peut supprimer des colonnes en perdant leurs données.
  \item \textbf{Un seul utilisateur PostgreSQL}, \code{nexora}, avec tous les droits sur tous les schémas : l'isolation entre services repose sur une convention, pas sur des permissions.
  \item \textbf{Aucune contrainte \code{CHECK}} : les bornes métier (\code{difficulty} entre 1 et 4, \code{selfGrade} entre 0 et 20, \code{masteryScore} entre 0 et 100, durées positives) ne sont contrôlées que par Zod dans le code.
  \item \textbf{URL complètes stockées} (\code{statementPdfUrl}, etc.) : elles deviendront fausses si l'adresse de MinIO change.
  \item \textbf{Écarts avec le frontend} : le frontend gère des conversations à plusieurs messages et des pièces jointes, alors que \code{contact_messages} ne prévoit qu'une seule \code{adminReply}. Les annales du frontend couvrent aussi Mines, Centrale, CCINP, la chimie et les langues, qui sont absents de \code{Epreuve} et de \code{cnc_exams}. La page des paramètres gère un avatar, une bio et des concours visés, qui ne sont pas stockés.
  \item \textbf{Pas de purge} des refresh tokens révoqués ou expirés.
  \item \textbf{\code{qcm_answers.choiceId}} n'a pas de clé étrangère vers \code{qcm_choices}.
  \item \textbf{Pas de sauvegarde} automatique.
\end{enumerate}

% -----------------------------------------------------------------------------
\subsection{Consignes pour les changements à venir}

\subsubsection{Les règles}

\begin{tcolorbox}[colback=soft,colframe=navy,title=\textbf{Règles à respecter pour toute modification de la base},fonttitle=\bfseries,breakable,arc=2pt]
\begin{enumerate}[label=\textbf{R\arabic*.},leftmargin=2.2em]
  \item \textbf{Un service, un schéma.} Chaque service ne modifie que son propre \code{schema.prisma}. Il ne lit ni n'écrit jamais, même en SQL, les tables d'un autre schéma : il passe par l'API ou par un événement.
  \item \textbf{Pas de clé étrangère entre schémas.} Une référence vers un autre service est un \code{String @db.Uuid} nommé \code{<entite>Id}, avec un index s'il sert à filtrer.
  \item \textbf{Toute modification passe par une migration} (\code{prisma migrate dev}). Les fichiers \code{prisma/migrations/*} sont \textbf{commités} avec le schéma. Une fois la migration initiale créée, \code{db push} est interdit, sauf sur une base locale jetable.
  \item \textbf{On ne modifie jamais une migration déjà appliquée} ou partagée : on en crée une nouvelle.
  \item \textbf{Relire le SQL avant de l'appliquer} (\code{--create-only}). Attention en particulier : renommer un champ dans \code{schema.prisma} produit un \code{DROP COLUMN} puis un \code{ADD COLUMN}, donc une \textbf{perte de données}. Il faut soit garder le nom en base avec \code{@map("ancienNom")}, soit remplacer ces deux lignes par un \code{ALTER TABLE ... RENAME COLUMN}.
  \item \textbf{Changements compatibles}, en trois temps (\emph{expand / contract}) : (1) ajouter la nouvelle colonne, nullable ou avec une valeur par défaut ; (2) déployer le code et recopier les données ; (3) dans une migration \emph{suivante}, rendre la colonne obligatoire et supprimer l'ancienne.
  \item \textbf{Énumérations} : ajouter une valeur est sans risque. En retirer ou en renommer une demande une migration écrite à la main, après avoir converti les lignes concernées. Toute modification d'une énumération dupliquée (\code{Filiere}, \code{SubjectCode}) se fait \textbf{dans les deux schémas et dans le frontend} (\code{lib/mappers.ts}).
  \item \textbf{Nommage} : modèle Prisma en \code{PascalCase} au singulier ; \code{@@map} en \code{snake_case} au pluriel ; champs en \code{camelCase} ; booléens en \code{is...} ou \code{has...} ; horodatages en \code{...At} ; jours calendaires en \code{...Date} avec \code{@db.Date}.
  \item \textbf{Index} : un index pour chaque clé étrangère et chaque colonne utilisée dans un \code{WHERE} ou un \code{ORDER BY} fréquent. On vérifie avec \code{EXPLAIN ANALYZE} sur un volume réaliste.
  \item \textbf{Contraintes métier en base} : ajouter les \code{CHECK} dans le SQL de la migration (voir ci-dessous) en plus de la validation Zod.
  \item \textbf{Fichiers} : jamais de PDF ni d'image en base. On stocke la \emph{clé} de l'objet MinIO (\code{statementPdfKey}) et on génère une URL signée au moment de la lecture.
  \item \textbf{Opérations sur plusieurs tables} dans une transaction, avec \code{prisma.$transaction} (par exemple : journal d'activité et Streak).
  \item \textbf{Données de départ} : chaque service a un \code{prisma/seed.ts} \emph{idempotent} (avec \code{upsert}), qu'on peut relancer sans créer de doublons. Les comptes de démonstration ne servent qu'en développement.
  \item \textbf{Sécurité} : aucun secret ni mot de passe en clair en base ou dans le dépôt. En production, chaque service a un utilisateur PostgreSQL limité à son propre schéma.
  \item \textbf{Sauvegarde} (\code{pg_dump}) avant toute migration qui supprime ou transforme des données.
\end{enumerate}
\end{tcolorbox}

\subsubsection{Procédure d'une modification}

\begin{figure}[H]\centering
\begin{tikzpicture}[node distance=4mm,
  st/.style={box,minimum width=2.6cm,minimum height=1.1cm,font=\scriptsize,text width=2.4cm}]
  \node[st,fill=royal!8] (1) {\textbf{1.} modifier\\\texttt{schema.prisma}};
  \node[st,fill=yellow!12,right=of 1] (2) {\textbf{2.} \texttt{migrate dev}\\\texttt{-{}-create-only}};
  \node[st,fill=warm!10,draw=warm,right=of 2] (3) {\textbf{3.} relire / compléter le SQL (RENAME, CHECK)};
  \node[st,fill=royal!8,right=of 3] (4) {\textbf{4.} \texttt{migrate dev}\\(applique)};
  \node[st,fill=royal!8,below=8mm of 4] (5) {\textbf{5.} \texttt{db:generate}\\code + tests};
  \node[st,fill=royal!8,left=of 5] (6) {\textbf{6.} vérifier l'écart\\(\texttt{migrate diff})};
  \node[st,fill=green!10,draw=okgreen,left=of 6] (7) {\textbf{7.} commit\\schéma + migrations};
  \node[st,fill=green!10,draw=okgreen,left=of 7] (8) {\textbf{8.} production :\\\texttt{migrate deploy}};
  \draw[arr] (1)--(2); \draw[arr] (2)--(3); \draw[arr] (3)--(4); \draw[arr] (4)--(5);
  \draw[arr] (5)--(6); \draw[arr] (6)--(7); \draw[arr] (7)--(8);
\end{tikzpicture}
\caption{Cycle de vie d'une modification du schéma}
\end{figure}

\subsubsection{Commandes de référence}
Toutes les commandes se lancent \textbf{dans le dossier du service} concerné, par exemple \code{backend/services/catalog-course-service}. Le préfixe \code{npx dotenv -e ../../.env --} charge le fichier \code{.env} commun.

\begin{lstlisting}[numbers=none,language={},caption={Étape préalable, une seule fois par service : partir de l'état actuel (baseline)}]
mkdir prisma\migrations\0_init
npx prisma migrate diff --from-empty --to-schema-datamodel prisma/schema.prisma --script -o prisma/migrations/0_init/migration.sql
npx dotenv -e ../../.env -- prisma migrate resolve --applied 0_init
# -> crée la table _prisma_migrations dans le schéma du service, sans toucher aux données
\end{lstlisting}

\begin{lstlisting}[numbers=none,language={},caption={Modification courante}]
npx dotenv -e ../../.env -- prisma migrate dev --create-only --name add_star_class_to_students
#   relire et éventuellement compléter prisma/migrations/<date>_add_star_class_to_students/migration.sql
npx dotenv -e ../../.env -- prisma migrate dev          # applique la migration et régénère le client
npx dotenv -e ../../.env -- prisma migrate diff --from-schema-datasource prisma/schema.prisma `
     --to-schema-datamodel prisma/schema.prisma --exit-code   # doit afficher : No difference detected
npx dotenv -e ../../.env -- prisma migrate deploy       # sur le serveur (CI / production)
\end{lstlisting}

\begin{lstlisting}[numbers=none,language=SQL,caption={Exemple de SQL ajouté à la main dans une migration : contraintes métier}]
ALTER TABLE "exercises"           ADD CONSTRAINT exercises_difficulty_chk   CHECK ("difficulty" BETWEEN 1 AND 4);
ALTER TABLE "written_submissions" ADD CONSTRAINT submissions_grade_chk      CHECK ("selfGrade" BETWEEN 0 AND 20);
ALTER TABLE "qcm_mastery"         ADD CONSTRAINT mastery_score_chk          CHECK ("masteryScore" BETWEEN 0 AND 100);
ALTER TABLE "cnc_exams"           ADD CONSTRAINT exams_duration_chk         CHECK ("durationMinutes" > 0);
\end{lstlisting}

\begin{lstlisting}[numbers=none,language={},caption={Consulter, sauvegarder et restaurer}]
docker exec -it backend-postgres-1 psql -U nexora -d nexora     # \dn  (schémas)   \dt auth.*  (tables)
npx dotenv -e ../../.env -- prisma studio                         # interface web de consultation
docker exec backend-postgres-1 pg_dump -U nexora -d nexora -Fc -f /tmp/nexora.dump
docker cp backend-postgres-1:/tmp/nexora.dump .\nexora_<date>.dump
docker exec backend-postgres-1 pg_restore -U nexora -d nexora --clean /tmp/nexora.dump
\end{lstlisting}

\subsubsection{Liste de contrôle avant de valider un changement}
\begin{itemize}[label=$\square$]
  \item La migration a été relue : aucun \code{DROP} non voulu, renommages traités, \code{CHECK} ajoutés.
  \item \code{npm run db:generate}, \code{npm run typecheck} et \code{npm test} passent.
  \item \code{prisma migrate diff ... --exit-code} indique « No difference detected ».
  \item Le \emph{seed} est mis à jour et peut être relancé sans créer de doublons.
  \item Les énumérations dupliquées et les correspondances du frontend sont synchronisées.
  \item Le schéma et le dossier \code{migrations/} sont commités \emph{ensemble}.
  \item Une sauvegarde existe si la migration supprime ou transforme des données.
\end{itemize}

\subsubsection{Changements de la base déjà identifiés}

\begin{longtable}{p{1.4cm} p{2.4cm} p{11.3cm}}
\toprule
\textbf{Quand} & \textbf{Schéma} & \textbf{Changement prévu}\\
\midrule\endhead
Avant S2 & tous & Migration initiale \code{0_init} (baseline) dans les 6 services ; suppression du compte de test ; tâche de purge des \code{refresh_tokens} révoqués ou expirés depuis plus de 30 jours.\\
S2 & \code{auth} & Décider pour l'« étoile » de filière : ajouter \code{starClass Boolean @default(false)} à \code{Student}. Ajouter les champs de profil gérés par le frontend : \code{avatarKey String?}, \code{bio String?}, \code{targetConcours String[]}.\\
S2 & \code{catalog} & \emph{Seed} des 4 matières avec leurs coefficients ; import des 26 chapitres, leçons et exercices du frontend ; \code{CHECK} sur \code{difficulty} ; remplacer les colonnes \code{...PdfUrl} par \code{...PdfKey} (expand/contract) ; index \code{(chapterId, orderIndex)} et \code{(lessonId)}.\\
S3 & \code{exam} & Arbitrer le périmètre des annales. Le CDC ne parle que du CNC, mais le frontend affiche aussi Mines, Centrale et CCINP. Si on les garde : ajouter une énumération \code{Concours} et une colonne \code{concours}, ainsi que les épreuves \code{CHIMIE} et \code{LANGUES}. Clé unique \code{(concours, year, filiere, epreuve, kind)} pour un import idempotent des 437 épreuves ; \code{CHECK} sur \code{durationMinutes} et \code{selfGrade}.\\
S4 & \code{contact} & Conversations à plusieurs messages : nouvelle table \code{ticket_messages} (\code{ticketId}, \code{authorId}, \code{authorRole}, \code{body}, \code{isInternalNote}, \code{createdAt}) et table \code{ticket_attachments} (\code{messageId}, \code{objectKey}, \code{fileName}, \code{size}). Recopier \code{adminReply} dans la nouvelle table, puis supprimer la colonne dans une migration ultérieure.\\
S4 & \code{analytics} & Aucun changement de structure nécessaire. Prévoir l'index \code{(userId, activityDate)} si la heatmap annuelle est lente.\\
S5 & \code{qcm} & Clé étrangère \code{qcm_answers.choiceId} vers \code{qcm_choices} ; \code{CHECK} sur \code{difficultyLevel} et \code{masteryScore} ; index \code{(sessionId)} sur \code{qcm_answers}.\\
À décider & — & Pochettes et pages personnelles du frontend : aucun service n'est prévu par le CDC. Soit nouveau schéma \code{workspace}, soit tables dans \code{auth}.\\
Production & tous & Un utilisateur PostgreSQL par service (\code{GRANT USAGE, CREATE ON SCHEMA catalog TO catalog_svc}), sauvegardes quotidiennes, \code{migrate deploy} lancé par la chaîne de déploiement.\\
\bottomrule
\caption{Évolutions de la base à prévoir, sprint par sprint}
\end{longtable}

\section{Ce qu'il reste à faire}

\subsection{Backend : sprints restants}

Le contrat d'API et les schémas existent déjà. Pour chaque sprint, le travail consiste à remplacer \code{notImplemented(...)} par de vrais handlers, en suivant le modèle de \code{auth-user-service} : \code{schemas.ts} (Zod), \code{*.service.ts} (règles métier), \code{routes.ts}.

\begin{longtable}{p{2cm} p{3cm} p{10.2cm}}
\toprule
\textbf{Sprint} & \textbf{Service} & \textbf{Travail restant}\\
\midrule\endhead
\textbf{S2}\newline sem. 3–5 & catalog-course & Lecture des matières, chapitres (filtres niveau et filière, \emph{sans} restriction d'accès), leçons, exercices double mode ; enregistrement d'une \code{ExerciseAttempt} qui publie un événement d'activité ; CRUD admin ; section Ouvrages (filtres matière, filière, niveau) ; script de \emph{seed} qui importe les 26 chapitres de \code{COURSES_DATABASE} du frontend ; envoi des PDF de cours vers MinIO (URL présignées).\\
\textbf{S3}\newline sem. 6–8 & exam-submission & Import des 437 épreuves depuis les manifestes JSON du frontend ; filtres année, filière, épreuve ; \code{ExamSession} dont la fin (\code{deadlineAt}) est calculée côté serveur (anti-triche) ; soumission LaTeX ou scan (MinIO) ; corrigé et barème renvoyés \emph{uniquement} après soumission ; auto-évaluation ; portefeuille « Mes Problèmes Rédigés » (UC07).\\
\textbf{S4}\newline sem. 9–10 & analytics-streak & \code{POST /analytics/activities} : \emph{upsert} de \code{DailyActivityLog} et appel à \code{updateStreak()} dans \textbf{une même transaction} ; heatmap annuelle, radar par matière (ventilé Sup/Spé), compteurs ; cache Redis du Dashboard (objectif : moins de 200\,ms au 95\textsuperscript{e} centile).\\
\textbf{S4} & admin-contact & Création de ticket avec \code{targetEntityRef} pré-rempli et l'une des 4 catégories ; cycle OPEN $\rightarrow$ IN\_PROGRESS $\rightarrow$ RESOLVED ; réponse de l'administrateur qui crée une \code{Notification} et envoie un e-mail (SMTP, par exemple Nodemailer).\\
\textbf{S5}\newline sem. 11–12 & adaptive-qcm & Session adaptative qui s'appuie sur \code{updateMastery()}, \code{nextDifficulty()} et \code{scheduleReview()} ; « QCM Flash Quotidien » (questions dont la révision est due) ; bilan \code{weakLessonsJson} ; événement Streak ; CRUD admin des questions.\\
\textbf{Transverse} & tous & Communication entre services par événements Redis (par exemple \code{activity.recorded}) plutôt que par appels HTTP directs ; migrations Prisma versionnées (\code{prisma migrate}) au lieu de \code{db push} ; journaux structurés ; tests d'intégration par service ; Dockerfile par service ; documentation OpenAPI.\\
\bottomrule
\caption{Travail backend restant, par sprint (planning du CDC, section 7.2)}
\end{longtable}

\subsubsection{Flux d'événements prévu pour le Streak}

Le CDC demande que \emph{toute} activité pédagogique valide le Streak (UC09). Plusieurs services produisent ces activités, mais seul \code{analytics-streak-service} doit tenir le compteur :

\begin{figure}[H]\centering
\begin{tikzpicture}[node distance=5mm]
  \node[skel,minimum width=3cm] (c) {catalog\\\scriptsize exercice résolu};
  \node[skel,minimum width=3cm,below=of c] (e) {exam\\\scriptsize copie soumise};
  \node[skel,minimum width=3cm,below=of e] (q) {qcm\\\scriptsize session terminée};
  \node[db,fill=red!6,right=16mm of e] (r) {Redis\\\scriptsize\texttt{activity.recorded}};
  \node[done,right=16mm of r,text width=3.6cm] (a) {analytics-streak\\\scriptsize upsert du journal\\+ \texttt{updateStreak()}\\(1 transaction)};
  \node[db,right=10mm of a,minimum width=1.6cm] (pg) {schéma\\analytics};
  \foreach \n in {c,e,q}{\draw[darr] (\n.east) -- (r.west);}
  \draw[darr] (r) -- (a); \draw[arr] (a) -- (pg);
\end{tikzpicture}
\caption{Publication et consommation des événements d'activité (à implémenter)}
\end{figure}

\subsection{Frontend : branchement sur l'API}

\subsubsection{Étape 1 : un client HTTP unique}

Il faut créer \code{frontend/lib/api.ts}. Ce module ajoute l'en-tête \code{Authorization}, appelle \code{/api/auth/refresh} une seule fois quand il reçoit un 401, puis rejoue la requête. Il remplacera les fonctions de \code{lib/auth.ts}. Le refresh token devrait idéalement être stocké dans un cookie \code{httpOnly} plutôt que dans \code{localStorage}.

\begin{lstlisting}[language=TypeScript,caption={Esquisse de \texttt{frontend/lib/api.ts} (à écrire)}]
const API = process.env.NEXT_PUBLIC_API_URL ?? 'http://localhost:4000/api'

export async function api<T>(path: string, init: RequestInit = {}, retry = true): Promise<T> {
  const res = await fetch(API + path, {
    ...init,
    headers: { 'Content-Type': 'application/json',
               Authorization: `Bearer ${getAccessToken()}`, ...init.headers },
  })
  if (res.status === 401 && retry && (await refreshTokens())) return api<T>(path, init, false)
  if (!res.ok) throw await res.json()          // { error, message, details? }
  return res.status === 204 ? (undefined as T) : res.json()
}
\end{lstlisting}

\subsubsection{Étape 2 : remplacer les données en dur}

\begin{longtable}{p{5.6cm} p{4.4cm} p{5cm}}
\toprule
\textbf{Emplacement dans le frontend} & \textbf{Donnée simulée} & \textbf{Endpoint cible}\\
\midrule\endhead
\code{lib/auth.ts}, \code{LoginPage.tsx}, \code{SignUpPage.tsx} & \code{DEMO_STUDENT}, faux jetons, mot de passe ignoré & \code{/api/auth/*} — \textbf{disponible}\\
\code{app/parametres/page.tsx} & profil dans \code{localStorage}, faux changement de mot de passe & \code{PATCH /api/users/me} — \textbf{disponible}\\
\code{app/admin/page.tsx} & aucune vérification du rôle & \code{GET /api/users/me} (rôle) — \textbf{disponible}\\
\code{app/page.tsx} (dashboard), \code{ActivityBar}, \code{StreakBadge}, \code{AchievementBadgesCard} & 72\,\%, 12 jours, 254\,h, heatmap pseudo-aléatoire & \code{/api/analytics/me/*}\\
\code{CourseExerciseQcmDoubleMode.tsx} : \code{COURSES_DATABASE}, \code{SUBJECTS} & 26 chapitres, 28 exercices & \code{/api/catalog/*}\\
même fichier : \code{exerciseEvaluations}, « Signaler » & auto-évaluation et erratum non conservés & \code{POST /api/catalog/}\allowbreak\code{exercises/:id/attempts}, \code{POST /api/contact/}\allowbreak\code{tickets}\\
même fichier : QCM flash et banque QCM & maîtrise +6/$-$4, chronomètre figé & \code{/api/qcm/*}\\
\code{lib/concours-annales-database.ts} et manifestes JSON & 437 épreuves, URL statiques & \code{/api/exams}\\
\code{ConcoursPdfExamViewer.tsx} & remise de copie et infractions en mémoire, candidat en dur & \code{/api/exams/:id/sessions}, \code{/submit}\\
\code{ExercisePdfTemplate.tsx}, \code{lib/classiques-problems.ts} & 10 problèmes, notes locales & \code{/api/exams?kind=PROBLEM}\\
\code{ContactAdminTemplate.tsx} & tickets en dur, réponse automatique au bout de 1,4\,s & \code{/api/contact/*}\\
\code{AdminBackOffice.tsx} & catalogue, annales, QCM et conversations en état local & routes admin des 4 services\\
\code{PochettesWorkspace.tsx} & dossiers et pages personnels & \emph{à arbitrer} (aucun service prévu par le CDC)\\
\bottomrule
\caption{Données simulées à remplacer par des appels à l'API}
\end{longtable}

\subsubsection{Étape 3 : harmoniser les énumérations}

\begin{table}[H]\centering\footnotesize
\begin{tabular}{l l l}
\toprule
\textbf{Champ} & \textbf{Frontend actuel} & \textbf{Backend (CDC)}\\
\midrule
Filière & \code{'MP*' | 'MP' | 'PSI*' | 'PSI' | 'TSI'} & \code{MP | PSI | TSI}\\
Année & \texttt{'1re année Sup' | '2e année Spé'} & \code{SUP | SPE}\\
Niveau d'un chapitre & \code{'sup' | 'spe'} & \code{SUP | SPE}\\
Matière & \code{maths, physique, si (ou sii), info} & \code{MATH | PHYS | SI | INFO}\\
Auto-évaluation & \code{mastered | with_help | review} & \code{MASTERED | WITH_HINT | TO_REVIEW}\\
Statut de ticket & \code{pending | in_progress | resolved} & \code{OPEN | IN_PROGRESS | RESOLVED}\\
\bottomrule
\end{tabular}
\caption{Correspondances à mettre en place (par exemple dans \texttt{frontend/lib/mappers.ts})}
\end{table}
Pour la filière, il faut décider si l'« étoile » (MP*) doit être conservée. Si oui, il faudrait ajouter un champ booléen \code{starClass} au modèle \code{Student}.

\subsection{Écarts restants avec le CDC}

\begin{table}[H]\centering\small
\begin{tabularx}{\linewidth}{X l l}
\toprule
\textbf{Exigence du CDC} & \textbf{Frontend} & \textbf{Backend}\\
\midrule
Inscription et connexion JWT + bcrypt (M1) & \partiel{} (simulé) & \fait\\
Basculeur « Tous les niveaux / Sup / Spé » & \partiel{} (pas de « Tous ») & \fait{} (\code{accessAllYears})\\
Organisation domaine $\rightarrow$ chapitre $\rightarrow$ leçon (M2) & \partiel{} & \fait{} (schéma)\\
Exercices en mode PDF avec énoncé et corrigé séparés & \afaire{} (planche HTML) & \fait{} (schéma)\\
Annales 2010–2026, filtre par filière (M3) & \partiel{} (MP seulement) & \afaire\\
« Mes Problèmes Rédigés » : soumission, archive, barème & \afaire & \afaire\\
Bibliothèque d'ouvrages (M4) & \afaire{} (absente) & \afaire{} (schéma prêt)\\
Record de Streak, ventilation Sup/Spé (M5) & \afaire & \afaire{} (algorithme prêt)\\
Signalement contextualisé et 4 catégories (M6) & \partiel & \afaire{} (schéma prêt)\\
QCM adaptatif et répétition espacée (P15) & \partiel{} & \afaire{} (algorithmes prêts)\\
Modules optionnels M7 (TIPE, simulateur de rang, oraux…) & \afaire & \afaire\\
\bottomrule
\end{tabularx}
\caption{Couverture du CDC par couche}
\end{table}

\subsection{Dette technique du frontend à traiter}

\begin{enumerate}
  \item \textbf{Découper les gros composants} : extraire les données de \code{CourseExerciseQcmDoubleMode.tsx} (5\,276 lignes) dans \code{lib/}, puis séparer les sous-vues (\code{CoursView}, \code{TdView}, \code{QcmBank}, \code{AnnalesPortal}). Faire de même pour \code{AdminBackOffice.tsx}.
  \item \textbf{Protéger les routes} : rediriger vers \code{/login} sans session ; contrôler le rôle ADMIN sur \code{/admin} (le backend refuse déjà, mais l'interface s'affiche quand même).
  \item \textbf{Déconnexion réelle} : appeler \code{POST /api/auth/logout}, puis \code{clearUserSession()}.
  \item \textbf{Supprimer le code mort} : \code{StudentDrawer}, \code{PersonalWorkspaceModal}, \code{glass-folder-cards}, \code{ui/button}, \code{@base-ui/react}.
  \item \textbf{Réactiver la vérification de types} : retirer \code{ignoreBuildErrors} de \code{next.config.mjs} et corriger les erreurs.
  \item \textbf{Corriger l'échappement LaTeX} : certaines formules contiennent \texttt{'\textbackslash forall'} au lieu de \texttt{'\textbackslash\textbackslash forall'} (en JavaScript, \texttt{\textbackslash f} devient un caractère de contrôle).
  \item \textbf{Déplacer les 570\,Mo de PDF} de \code{public/} vers MinIO, servis par \code{exam-submission-service} avec des URL signées.
\end{enumerate}

\subsection{Planning proposé}

\begin{figure}[H]\centering
\begin{tikzpicture}[x=1.15cm,y=0.62cm]
  \foreach \w in {1,...,12}{\node[font=\scriptsize] at (\w-0.5,0.8) {S\w}; \draw[black!12] (\w,0.4) -- (\w,-7.6);}
  \draw[black!12] (0,0.4) -- (0,-7.6);
  \newcommand{\ganttbar}[5]{\fill[#4,rounded corners=1.5pt] (#1,-#3+0.32) rectangle (#2,-#3-0.32); \node[font=\scriptsize,anchor=east] at (-0.1,-#3) {#5};}
  \ganttbar{0}{2}{0.5}{okgreen!70}{Sprint 1 : gateway + auth}
  \ganttbar{2}{5}{1.5}{royal!60}{Sprint 2 : catalog}
  \ganttbar{5}{8}{2.5}{royal!60}{Sprint 3 : exam}
  \ganttbar{8}{10}{3.5}{royal!60}{Sprint 4 : analytics + contact}
  \ganttbar{10}{12}{4.5}{purple!50}{Sprint 5 : QCM + options}
  \ganttbar{2}{10}{5.5}{warm!50}{Branchement du frontend}
  \ganttbar{4}{12}{6.5}{black!25}{Tests, Docker, documentation}
  \draw[warm,line width=1pt] (2,0.5) -- (2,-7.4) node[below,font=\scriptsize]{aujourd'hui};
\end{tikzpicture}
\caption{Avancement par rapport au planning du CDC (le Sprint 1 est terminé)}
\end{figure}

\section{Guide de démarrage}

\subsection{Prérequis}
Node.js 24 ou plus récent, npm 11, et Docker Desktop (pour PostgreSQL, Redis et MinIO).

\subsection{Backend}
\begin{lstlisting}[numbers=none,language={}]
cd backend
copy .env.example .env      # puis remplir JWT_ACCESS_SECRET et INTERNAL_SECRET
                            # node -e "console.log(require('crypto').randomBytes(48).toString('hex'))"
npm install
npm run infra:up            # PostgreSQL :5434, Redis :6379, MinIO :9000 (console :9001)
npm run db:generate
npm run db:push
npm run db:seed             # admin@learning-cnc.ma / AdminPassword2026!
                            # achraf.lemrani@cpge-cnc.ma / StudentPassword2026!
npm run dev                 # gateway :4000 + 6 services :4001..4006
\end{lstlisting}

Vérification rapide :
\begin{lstlisting}[numbers=none,language={}]
curl http://localhost:4000/api/health
curl -X POST http://localhost:4000/api/auth/login -H "Content-Type: application/json" \
     -d "{\"email\":\"admin@learning-cnc.ma\",\"password\":\"AdminPassword2026!\"}"
\end{lstlisting}

\subsection{Frontend}
\begin{lstlisting}[numbers=none,language={}]
cd frontend
npm install                 # postinstall : copie les fichiers de pdf.js dans public/pdfjs
npm run dev                 # http://localhost:3000 (ou 3001 si le port 3000 est pris)
\end{lstlisting}

\subsection{Ajouter une route dans un service (méthode)}
\begin{enumerate}
  \item Écrire le schéma Zod de l'entrée dans \code{src/schemas.ts}.
  \item Écrire la fonction métier dans \code{src/<domaine>.service.ts}. Elle utilise \code{prisma} et lève une \code{HttpError} en cas de règle violée.
  \item Dans \code{src/routes.ts} (ou \code{index.ts}), remplacer \code{todo} par :
\begin{lstlisting}[language=TypeScript,numbers=none]
router.get('/catalog/chapters', async (req, res) => {
  res.json(await catalog.listChapters(chapterQuerySchema.parse(req.query)))
})
\end{lstlisting}
  \item Si le schéma Prisma a changé : \code{npm run db:generate} puis \code{npm run db:push}.
  \item Ajouter un test (\code{*.test.ts}), puis lancer \code{npm run typecheck} et \code{npm test}.
\end{enumerate}

\begin{note}[Pour relire le code dans le bon ordre]
\begin{enumerate}
  \item \code{backend/gateway/src/index.ts}, puis \code{authenticate.ts} : comment une requête entre.
  \item \code{backend/packages/common/src/} : le socle de tous les services.
  \item \code{backend/services/auth-user-service/src/routes.ts}, puis \code{auth.service.ts} et \code{tokens.ts} : un service complet.
  \item \code{prisma/schema.prisma} de chaque service : le modèle de données du CDC.
  \item \code{streak.ts} et \code{adaptive-engine.ts}, avec leurs tests : les algorithmes.
  \item \code{frontend/app/page.tsx}, puis \code{components/CourseExerciseQcmDoubleMode.tsx} : l'interface.
\end{enumerate}
\end{note}


\end{document}
